\documentclass[11pt]{article}
\usepackage{fullpage}
\usepackage{amsmath,amssymb,amsthm,amsfonts}
\usepackage{mathtools}
\usepackage{hyperref}
\hypersetup{hidelinks}
\allowdisplaybreaks
\theoremstyle{plain}
\newtheorem{theorem}{Theorem}[section]
\newtheorem{proposition}[theorem]{Proposition}
\newtheorem{lemma}[theorem]{Lemma}
\newtheorem{corollary}[theorem]{Corollary}
\newtheorem{conjecture}[theorem]{Conjecture}
\newtheorem{problem}[theorem]{Open Problem}
\theoremstyle{definition}
\newtheorem{definition}[theorem]{Definition}
\newtheorem{example}[theorem]{Example}
\theoremstyle{remark}
\newtheorem{remark}[theorem]{Remark}
\newcommand{\C}{\mathbb{C}}
\newcommand{\Q}{\mathbb{Q}}
\newcommand{\A}{\mathbb{A}}

\title{Counterexamples to the Jacobian conjecture\\ in dimensions greater than two}
\author{Shuhong Gao\thanks{AI disclosure: The main idea and framework are
due to the author, and Claude Fable 5 assisted in the proofs and in the writing
up of the paper.}\\[4pt]
\normalsize School of Mathematical and Statistical Sciences\\
Clemson University, Clemson, SC 29634-0975\\
Email: \texttt{sgao@clemson.edu}}
\date{July 31, 2026}

\begin{document}
\maketitle

\begin{abstract}
The Jacobian conjecture, open since 1939, asks whether every polynomial map of
$\C^n$ whose Jacobian determinant is a nonzero constant must have a polynomial
inverse. It was refuted in dimension three by Alp\"oge on July 19, 2026, with
an infinite family by Gallagher (July 20) and a geometric explanation by
Speyer (July 23): the counterexample sweeps the tangent lines of a plane
curve --- a map that classical duality forces to hit most points several
times. We give a self-contained account of this tangent-sweep mechanism and
generalize it from plane curves to direction fields on hypersurfaces. The
resulting construction produces counterexamples in every dimension greater
than two and, in each dimension, of arbitrarily large geometric degree (the
number of preimages of a typical point). We work it out in five new explicit
maps: one three-dimensional of degree four, two four-dimensional of degrees
five and ten, and two five-dimensional of degrees six and twelve. The
counterexamples provide explicit examples of \'etale coverings
$\C^n\to\C^n$ that are not proper: they are everywhere unramified, and fail
to be injective only through points escaping to infinity. All identities were
verified in exact rational arithmetic; an appendix determines exact fiber
structures through Gr\"obner bases.
\end{abstract}

\section{Introduction}

Let $F\colon\C^n\to\C^n$ be a polynomial map. Throughout this paper we use the
following conventions. The value $F(x)$ is written as a \emph{column} vector of
the $n$ components of $F$; consequently each partial derivative
$\partial_iF:=\partial F/\partial x_i$ is again a column vector, and the Jacobian
matrix of $F$ is
\[
J(F)\;=\;\bigl(\partial_1F,\;\partial_2F,\;\dots,\;\partial_nF\bigr),
\]
the square matrix whose columns are the partial derivatives, with the variables
taken in the order in which they are displayed. More generally, for column
vectors $V_1,\dots,V_k$ of equal length, $(V_1,\dots,V_k)$ denotes the matrix
with these columns; adjoining a column $\Delta$ to a matrix $J$ is written
$(\Delta,J)$. If $F$ has a polynomial inverse, the chain rule forces $\det J(F)$
to be a nonzero constant. The \emph{Jacobian conjecture}, raised by Keller
\cite{Keller1939}, asserts the converse: a polynomial map with
$\det J(F)\in\C^\times$ (a \emph{Keller map}) is a polynomial automorphism. The
problem became widely known through the survey of Bass, Connell and Wright
\cite{BCW1982} and the book of van den Essen \cite{Essen2000}.

On July 19, 2026, the conjecture was refuted in dimension three by an explicit
counterexample announced by L.~Alp\"oge \cite{Alpoge2026}. The next day,
Gallagher \cite{Gallagher2026} gave a uniform one-variable construction
producing, for every integer $d\ge3$, an explicit counterexample of geometric
degree $d$; expositions by Tao (July 21) \cite{Tao2026} and Speyer
(July 23) \cite{Speyer2026} followed. The first
counterexample is a polynomial map $F\colon\C^3\to\C^3$ with
$\det J(F)\equiv-2$ that identifies three distinct points. Its mechanism, isolated
and named the \emph{tangent sweep} by Speyer \cite{Speyer2026}, is strikingly
classical: one sweeps the tangent lines of a plane curve --- a map that
is unavoidably many-to-one, by projective duality --- and then conjugates by
monomial maps so that the Jacobian factor picked up by the sweep is cancelled
exactly. The price of the cancellation is that the ramification of the sweep is
pushed to infinity: the resulting Keller map is everywhere unramified but not
proper, and injectivity fails through escape to infinity, a loophole the Jacobian
hypothesis never excluded.

This paper has several goals. First (\S\ref{sec:sweep}), we give
a self-contained account of the tangent-sweep mechanism in the plane and of the
first counterexample (\S\ref{sec:first}), including a complete description of its
fibers: the generic
fiber has $3$ points, the fiber degenerates to $1$ point over an irreducible
surface $\Sigma$ in the target and to the empty set over a rational curve
$\mathcal{C}\subset\overline{\Sigma}$, and no fiber has exactly $2$ points. The
entire stratification is the projective duality of a cuspidal cubic read through a
monomial twist. Second (\S\ref{sec:deg4}), we construct along the same lines a new
three-dimensional Keller map of geometric degree \emph{four}, based on a rational
quartic curve with two cusps and a node; throughout the paper, the
\emph{geometric degree} of a generically finite polynomial map is the number of
points in its generic fiber, so ``geometric degree $d$'' and ``generic fiber of
exactly $d$ points'' are synonymous. Third (\S\ref{sec:dim4}), we give a general construction, valid in every
dimension $n>2$, based on sweeping tangent direction fields on parametrized
hypersurfaces. The methods are developed in
\S\ref{subsec:sweephyp}--\S\ref{subsec:stage}: prescribed direction fields and
their tangency criterion, formulated as exterior-product identities for any
direction field with unimodular component row, which split the construction
into escape branches (\S\ref{subsec:forms}); and the side conditions of the
monomial twist, organized into the discrete data and the stage equations
(\S\ref{subsec:stage}). Concrete examples follow in dimension four
(\S\ref{subsec:n4}): sweeping a tangent direction field on a suitable
\emph{non-cylindrical ruled surface} in $\C^3$ produces a Keller map
$F_4\colon\C^4\to\C^4$ with $\det J(F_4)\equiv-\tfrac{44}{9}$ and geometric
degree \emph{five} --- since geometric degree is invariant under composition
with polynomial automorphisms, $F_4$ is not equivalent to any product
$\Phi\times\mathrm{id}$ of a three-dimensional Keller map $\Phi$ of geometric
degree $\neq5$ with the identity, in particular not to the trivial extension
of Alp\"oge's example --- while specializing the second ($M$-)branch to the
direction field $(1,w_1,w_2)^{\mathsf T}$ yields a further four-dimensional
counterexample $F_5$ of geometric degree \emph{ten} (\S\ref{subsec:specII}).
The framework is then carried into dimension five (\S\ref{subsec:n5}), where a
new phenomenon appears: the direction field decides not only \emph{which}
escape branch is available but whether a branch is \emph{rigid}. For the mixed
field $\Delta=(1,w_1,w_1^2,w_2)^{\mathsf T}$ --- the first direction field that is
neither curve-like nor tautological --- the top branch vanishes identically,
and we prove that on the middle branch the governing equation degenerates,
on a natural characteristic-invariance locus, to a linear transport equation all
of whose solutions are suspensions of four-dimensional sweeps --- whether the
entire middle branch is rigid is left as an open problem
(Problem~\ref{prob:rigid5}); the bottom
branch, by contrast, is flexible, and we construct from it an explicit
five-dimensional Keller map $F_6$ with $\det J(F_6)\equiv-290$, of geometric
degree six, whose fiber over a punctured coordinate axis has exactly four
points. Finally (\S\ref{subsec:curvefam}), the curve-type
direction field $(1,w_1,\dots,w_1^{n-2})^{\mathsf T}$ admits a \emph{reductive}
construction, uniform in every dimension: its normalization is
\emph{equivalent} to prescribing a
$w_1$-family of constant-Jacobian maps of $\C^{n-3}$ plus one free potential,
its fibers are governed by a single univariate polynomial generalizing the
tangency polynomial of the plane theory, and it yields a second
five-dimensional counterexample $F_7$, of geometric degree twelve; since the
transverse data need only have constant Jacobian, Keller counterexamples in
$n-3$ variables can themselves serve as data, so the construction propagates
its own counterexamples along $n\mapsto n+3$. In particular, the construction
introduced here is not limited to the five explicit maps of this paper: it
produces counterexamples in every dimension $n>2$ and, within each fixed
dimension, of arbitrarily large geometric degree --- for $n=3$ through
Gallagher's family, and for every $n\ge4$ through the curve-type field, whose
free data may be chosen of arbitrarily high degree, the geometric degree being
the degree of the univariate fiber polynomial of
Proposition~\ref{prop:curvefib}.

All polynomial identities asserted below (Jacobian determinants, divisibilities,
fiber counts at rational points) were verified by exact rational arithmetic; the
verification scripts are available from the author. Appendix~\ref{app:fibers}
defines graph ideals and analyzes exact fiber sizes: for $F$ we give the full
(six-element) reduced lexicographic Gr\"obner basis certifying the stratification
of \S\ref{sec:first}; for $G$ we prove a complete fiber-size theorem
(Theorem~\ref{thm:fibG}) showing that $G$ attains \emph{every} fiber size in
$\{4,3,2,1,0\}$; for $F_4$--$F_7$, whose Gr\"obner bases are out of reach, we
summarize the univariate elimination data behind the generic counts, the
complete stratifications being deferred to a subsequent version.

\section{Background and previous work}\label{sec:background}

We recall what was known before 2026; see \cite{Essen2000} for a comprehensive
treatment. Wang \cite{Wang1980} proved the conjecture for maps of degree at most
$2$ in any dimension. Bass--Connell--Wright \cite{BCW1982} and Yagzhev
\cite{Yagzhev1980} showed that it suffices to prove the conjecture for maps of the
form $x+H(x)$ with $H$ cubic homogeneous (in all dimensions simultaneously), and
Dru\.zkowski \cite{Druzkowski1983} reduced further to the case where each
component of $H$ is the cube of a linear form. In dimension two, Moh
\cite{Moh1983} verified the conjecture for maps of degree at most $100$; the
two-dimensional case remains open and is untouched by the counterexamples
discussed here, which exist only in dimension $\ge3$ (by Wang's theorem, degree
$2$ examples are impossible, and the known constructions produce degree $\ge3$).

Two structural facts frame the counterexample. A theorem of
Bia{\l}ynicki-Birula and Rosenlicht \cite{BBR1962} states that an injective
polynomial map $\C^n\to\C^n$ is bijective with polynomial inverse; hence a Keller
counterexample must fail injectivity. Since a Keller map is \'etale (unramified),
it is a local biholomorphism everywhere, so non-injectivity can only occur through
failure of \emph{properness}: distinct preimages cannot collide over any target
point, they can only escape to infinity. The locus of points over which a
polynomial map is not proper was studied by Jelonek \cite{Jelonek1993}, who showed
it is either empty or a hypersurface. In the real category, Pinchuk
\cite{Pinchuk1994} had already constructed a non-injective polynomial map of
$\mathbb{R}^2$ with everywhere nonvanishing (nonconstant) Jacobian, an early
warning that unramifiedness does not force injectivity. The counterexamples below
realize this scenario over $\C$ with \emph{constant} Jacobian: they are \'etale of
geometric degree $d>1$ onto their (open, dense) image, with the fiber count
dropping only across the Jelonek hypersurface.

\section{The tangent sweep of a plane curve}\label{sec:sweep}

In this section we describe the main idea behind Alp\"oge's counterexample and
Gallagher's family \cite{Gallagher2026}, following the \emph{tangent sweep}
picture due to Speyer \cite{Speyer2026}.

\subsection{The sweep and its Jacobian}
Let $p,q\in\C[w]$ with $p$ nonconstant, and let
$\mathcal{K}\colon w\mapsto(p(w),q(w))^{\mathsf T}$ be a parametrized plane curve.
Throughout this section we assume
\[
\deg q\;\le\;\deg p+1 ,
\]
so that the tangent slope $q'/p'$ stays bounded as $w\to\infty$: the tangent
direction of $\mathcal K$ does not degenerate to the vertical at its place at
infinity. The assumption is exactly the compatibility condition for the normal
form below: comparing degrees in the normalization $q'=\tfrac{w}{2}p'$ shows
that any pair satisfying it has $\deg q=\deg p+1$ \emph{exactly}, while a pair
with $\deg q>\deg p+1$ cannot be brought to that form; and it is what makes
the degree and leading-coefficient statements of Lemma~\ref{lem:fibroots}
--- hence the fiber counts --- uniform over all targets.

The \emph{tangent sweep} of $\mathcal{K}$ is the map
sending (point of $\mathcal K$, position along its tangent line) to the plane:
\[
T_0(\gamma,w)=\bigl(p(w)+\gamma p'(w),\;q(w)+\gamma q'(w)\bigr)^{\mathsf T},
\qquad
\det J(T_0)=\gamma\,(p'q''-p''q') ,
\]
the variables being ordered $(\gamma,w)$.
The factor $p'q''-p''q'$ vanishes at inflection points. Normalizing the direction
field to $\delta=(2,w)^{\mathsf T}$, which is tangent exactly when $q'=\tfrac{w}{2}p'$, i.e.
\begin{equation}\label{eq:normalization}
q(w)=\int_0^w \tfrac{s}{2}\,p'(s)\,ds ,
\end{equation}
gives the inflection-free normal form
\[
S(\gamma,w)=\bigl(p(w)+2\gamma,\;q(w)+\gamma w\bigr)^{\mathsf T},\qquad \det J(S)=2\gamma .
\]
Thus the sweep is unramified except over $\gamma=0$, which maps onto the curve
itself, and its Jacobian is a \emph{coordinate}.

\subsection{Fibers and duality}\label{subsec:fibdual}

Fix a target point $(X,Y)\in\C^2$ and solve $S(\gamma,w)=(X,Y)$. The first
coordinate determines the position along the tangent line,
$\gamma=\tfrac12(X-p(w))$; substituting into the second coordinate leaves a
single equation in $w$, the \emph{tangency polynomial}
\begin{equation}\label{eq:tangency}
W_{X,Y}(w)\;=\;q(w)+\tfrac{w}{2}\bigl(X-p(w)\bigr)-Y .
\end{equation}
Its meaning is geometric: the tangent line of $\mathcal K$ at the parameter
value $w$ is the line $\{\mathcal K(w)+\gamma\,(2,w)^{\mathsf T}:\gamma\in\C\}$, and
$W_{X,Y}(w)=0$ says exactly that this line passes through $(X,Y)$. A root
$w_0$ of $W_{X,Y}$ is thus a \emph{parameter of tangency}: it names the point
$\mathcal K(w_0)$ of the curve at which some tangent line through $(X,Y)$
touches $\mathcal K$. (Picture a point outside a parabola: two tangent lines
pass through it, touching the parabola at two points; the two parameter values
naming those points of contact are the two roots of $W_{X,Y}$.)

\begin{lemma}\label{lem:fibroots}
For every $(X,Y)\in\C^2$ the assignment
$w_0\mapsto\bigl(\tfrac12(X-p(w_0)),\,w_0\bigr)$ is a bijection from the set
of distinct roots of $W_{X,Y}$ onto the fiber $S^{-1}(X,Y)$. If $\deg p=d$
then $\deg_wW_{X,Y}=d+1$ with leading coefficient
$-\mathrm{lc}(p)/(2(d+1))$, independent of $(X,Y)$: every fiber of $S$ is
finite of cardinality at most $d+1$, and $S$ is generically $(d+1)$-to-one.
Moreover the normalization \eqref{eq:normalization} gives the exact derivative
identity
\begin{equation}\label{eq:Wprime}
W_{X,Y}'(w)\;=\;\tfrac12\bigl(X-p(w)\bigr)\;=\;\gamma ,
\qquad
W_{X,Y}''(w)\;=\;-\tfrac12\,p'(w) .
\end{equation}
\end{lemma}

The identity \eqref{eq:Wprime} --- the derivative of the tangency polynomial
\emph{is} the sweep multiplier --- organizes all fiber degenerations, and is
used repeatedly later:
\begin{itemize}
\item \emph{Multiple roots lie on the curve.} $w_0$ is a multiple root iff
$W=W'=0$ there, i.e.\ iff $\gamma=0$ and $(X,Y)=\mathcal K(w_0)$: the
discriminant locus of $W_{X,Y}$ in the $(X,Y)$-plane is the swept curve
itself (the envelope of its tangent lines). Concretely, for the cuspidal cubic
of \S\ref{sec:first}, where $W_{X,Y}$ is proportional to $w^3-2w^2+Xw-2Y$,
one has the exact identity
$\operatorname{disc}_w\bigl(w^3-2w^2+Xw-2Y\bigr)=-4\,E(X,Y)$ with $E$ the
implicit equation of the cubic --- the sweep-level source of the discriminant
factorizations in Theorem~\ref{thm:strat3} and Appendix~\ref{app:fibers}.
\item \emph{Contact bookkeeping.} A point of $\mathcal K$ sees its own
tangent doubly: at a smooth point $\mathcal K(w_0)$ the root $w_0$ of
$W_{\mathcal K(w_0)}$ has multiplicity exactly $2$ (by \eqref{eq:Wprime},
$W''(w_0)=-\tfrac12p'(w_0)\neq0$). At a cusp ($p'(w_0)=0$, hence also
$q'(w_0)=0$) the multiplicity is $\ge3$: the cusp sees its cuspidal tangent
triply. At a node, \emph{two} parameters $w_1\neq w_2$ are each double roots.
\item \emph{The multiplier detects the point of contact.} A root has
$\gamma=0$ iff the target $(X,Y)$ is the point of tangency itself; these are
exactly the solutions lying on the ramification locus of $S$.
\end{itemize}

Counting distinct roots gives the fiber sizes of the sweep, and --- after the
twist of \S\ref{subsec:twist}, which deletes precisely the solutions with
$\gamma=0$ --- of the Keller maps built from it:

\begin{proposition}[fiber sizes]\label{prop:fibsizes}
Let $\deg p=d$. The fiber of the sweep $S$ over $(X,Y)$ has cardinality
\[
d+1\ \ \text{if }(X,Y)\notin\mathcal K,\qquad
d\ \ \text{at a smooth point of }\mathcal K,\qquad
d-1\ \ \text{at a cusp or a node},
\]
each contact of order $k$ collapsing $k$ simple roots into one multiple root.
The solutions with $\gamma=0$ are exactly the multiple roots, so for the
twisted map the corresponding sheets escape to infinity and the fiber sizes
become
\[
d+1\ \ (\text{generic}),\qquad
d-1\ \ (\text{smooth curve points}),\qquad
d-2\ \ (\text{cusps}),\qquad
d-3\ \ (\text{nodes}),
\]
away from the degeneration loci of the monomial conjugation itself, which can
contribute further strata. Both refinements are visible in the examples: for
$d=2$ (Theorem~\ref{thm:strat3}) the sizes are $3,1,0$ --- no nodes, and no
extra strata; for $d=3$ (Theorem~\ref{thm:fibG}) they are $4,2,1,0$,
\emph{plus} a twist stratum of size $3$ over a hyperplane invisible at sweep
level.
\end{proposition}

Behind the generic count stands classical duality: $d+1$ is the \emph{class}
of $\mathcal{K}$, the degree of the dual curve $\mathcal{K}^\vee$, computable
from the Pl\"ucker formula $m=\deg(\deg-1)-2\delta-3\kappa$ for a curve with
$\delta$ nodes and $\kappa$ cusps; the constancy of the leading coefficient of
$W_{X,Y}$ says that no tangency parameter escapes to $w=\infty$ as the target
moves, so the count $d+1$ holds over \emph{every} target, with multiplicity.

\subsection{The monomial twist}\label{subsec:twist}
The sweep is not a Keller map ($\det J(S)=2\gamma$), but the factor $\gamma$ can be
cancelled by conjugation with monomial maps. Pad the sweep to three variables and
compose:
\begin{multline*}
\begin{pmatrix}x\\ y\\ z\end{pmatrix}
\;\xrightarrow{\ \text{monomial}\ }\;
\begin{pmatrix}x\\ xy\\ x^2z\end{pmatrix}
\;\xrightarrow{\ \text{affine}\ }\;
\begin{pmatrix}x\\ \gamma\\ u\end{pmatrix}
\;\xrightarrow{\ w=\gamma u\ }\;
\begin{pmatrix}x\\ \gamma\\ w\end{pmatrix}\\
\xrightarrow{\ \text{padded sweep}\ }\;
\begin{pmatrix}\gamma x\\ P\\ Q\end{pmatrix}
\;\xrightarrow{\ \text{twist}\ }\;
\begin{pmatrix}C\\ P/C\\ Q/C^2\end{pmatrix},
\end{multline*}
where $\gamma=\gamma_0+a\,xy+b\,x^2z$, $u=1+xy$, $P=p(w)+2\gamma$,
$Q=q(w)+\gamma w$ and $C=\gamma x$. The Jacobians multiply to
\[
C^{-3}\cdot 2\gamma^2\cdot\gamma\cdot x^3\cdot(-b)\;=\;-2b,
\]
a nonzero constant. (Permuting the components of the final map changes the sign
of the determinant by the sign of the permutation; this accounts for
$\det J(F)=-2$ below, where the traditional component order reverses the one
displayed here.) The twist requires $C\mid P$ and $C^2\mid Q$; these
divisibilities amount to finitely many \emph{linear} conditions on the
coefficients of $p$ (and on $\gamma_0,a$), the ``side conditions'' of
\cite{Speyer2026}; \S\ref{subsec:stage} systematizes these conditions in every
dimension. Because the
ramification locus $\gamma=0$ of the sweep is exactly where the twisted coordinate
$C=\gamma x$ degenerates, the composite is everywhere unramified; the
$d+1$ sheets of the sweep survive, but sheets that would have merged over the
curve now escape to infinity. \emph{The construction converts ramification into
non-properness.} Varying $d=\deg p$ produces Keller counterexamples of geometric
degree $d+1$ for every $d\ge2$ --- this is Gallagher's family
\cite{Gallagher2026} --- which defeats all approaches to the conjecture through
degree bounds.

\subsection{The first counterexample and its fiber structure}\label{sec:first}

We first run the scheme of \S\ref{subsec:twist} with $\deg p=2$, which recovers
the counterexample of Alp\"oge \cite{Alpoge2026}. The solution of the side
conditions in closed form is
\[
p(w)=-3w^2+4w,\qquad
q(w)=-w^3+w^2,\qquad
\gamma=2-3xy-x^2z,\quad u=1+xy,\quad w=\gamma u
\]
($q$ is \eqref{eq:normalization} for this $p$), and the map, in Alp\"oge's
original component order,
\[
F\;=\;\Bigl(\frac{q(w)+\gamma w}{(\gamma x)^2},\;\;
\frac{p(w)+2\gamma}{\gamma x},\;\;\gamma x\Bigr)^{\mathsf T}\colon\ \C^3\longrightarrow\C^3
\]
(the sweep order of \S\ref{subsec:twist} lists the components in reverse; the
reversal is a transposition of the outer components and multiplies the Jacobian
determinant by $-1$).

\begin{theorem}\label{thm:alpoge}
$F$ is a polynomial map with components of degrees $7,6,4$; its Jacobian
determinant is identically $-2$; and its generic fiber consists of exactly
$3$ points.
\end{theorem}

A few comments on the ingredients. In the sweep order the determinant is
$-2b=2$ with $b=-1$, and the component reversal contributes the sign $-1$.
The divisibilities required by the twist are explicit:
$p(w)+2\gamma=\gamma\bigl(-3\gamma u^2+4u+2\bigr)$ and
$q(w)+\gamma w=\gamma^2\bigl(-\gamma u^3+u^2+u\bigr)$, and both bracketed
factors vanish at $x=0$. The tangency
polynomial \eqref{eq:tangency} is, up to a constant factor, the cubic
\[
w^3-2w^2+Xw-2Y ,
\]
whose leading coefficient is constant, so every target point admits three
tangency parameters counted with multiplicity.

Written out in full --- matching the form in \cite{Speyer2026,Tao2026} --- the
map is $F=(f_1,f_2,f_3)^{\mathsf T}$ with $f_3=\gamma x$,
$f_2f_3=p(w)+2\gamma$, $f_1f_3^{\,2}=q(w)+\gamma w$:
\begin{align*}
f_1 &= (1+xy)^3 z + y^2(1+xy)(4+3xy),\\
f_2 &= y + 3x(1+xy)^2 z + 3xy^2(4+3xy),\\
f_3 &= 2x - 3x^2y - x^3 z ,
\end{align*}
and
$F(0,0,-\tfrac14)=F(1,-\tfrac32,\tfrac{13}{2})=F(-1,\tfrac32,\tfrac{13}{2})
=(-\tfrac14,0,0)$: three distinct points with a common image. The underlying
curve $\mathcal{K}=(p,q)^{\mathsf T}$ is a rational \emph{cuspidal cubic} with one cusp, at
$(p,q)(\tfrac23)=(\tfrac43,\tfrac4{27})$, and no node; the Pl\"ucker formula
gives class $m=3\cdot2-3\cdot1=3$, matching the fiber count.

The fibers of $F$ are completely described by the following stratification, which
we established by exact computation (a lexicographic Gr\"obner basis of the graph
ideal, together with symbolic lifting certificates over the function fields of the
strata).

\begin{theorem}\label{thm:strat3}
Let $c_3(v)=27v_1^2v_3^2-18v_1v_2v_3+16v_1+v_2^3v_3-v_2^2$ and
$a_2(v)=54v_1v_3^2-18v_2v_3+16=\partial c_3/\partial v_1$, and let
$\sigma(v)=(v_2v_3,\,v_1v_3^2)^{\mathsf T}$. Then:
\begin{enumerate}
\item $\operatorname{disc}_x$ of the elimination cubic of the graph ideal equals
$-c_3a_2^2$, and $a_2^2-4(4-3v_2v_3)^3=108v_3^2c_3$;
\item $v_3^2c_3=E\circ\sigma$ where $E(X,Y)=X^3-X^2-18XY+16Y+27Y^2$ is the
implicit equation of $\mathcal{K}$, and $a_2$ is (a constant multiple of) the
equation of the cuspidal tangent line pulled back by $\sigma$;
\item the fiber of $F$ over $v$ has exactly $3$ points if $c_3(v)\ne0$ (with two
of them sharing their $x$-coordinate precisely when $a_2(v)=0$), exactly $1$
point if $c_3(v)=0$ and $v\notin\mathcal{C}$, and is empty if
$v\in\mathcal{C}=\{(\tfrac4{27}t^{-2},\tfrac43t^{-1},t):t\ne0\}$, the
$\sigma$-preimage of the cusp. In particular no fiber has exactly $2$ points, the
image of $F$ is $\C^3\setminus\mathcal{C}$, and the Jelonek non-properness set is
the irreducible surface $\{c_3=0\}=\sigma^{-1}(\mathcal{K})$.
\end{enumerate}
\end{theorem}

The mechanism is the contact geometry of
Proposition~\ref{prop:fibsizes} (here $d=2$: sizes $3,1,0$): crossing
$\sigma^{-1}(\mathcal K)$, a double contact makes two sheets escape together
(whence the exact square $a_2^2$ in the discriminant, and the impossibility of
fiber size $2$); over the cusp the contact is triple and all three sheets escape.

\subsection{A three-dimensional example of geometric degree four}\label{sec:deg4}

We now run the scheme of \S\ref{subsec:twist} with $\deg p=3$. Gallagher's
family \cite{Gallagher2026} contains examples of every geometric degree
$n\ge3$, including a degree-four instance; the example below was obtained
independently by solving the side conditions directly, and is recorded here with
its verification data. The solution in closed form is:
\[
p(w)=w^3-6w^2+6w,\qquad
q(w)=\tfrac38 w^4-2w^3+\tfrac32 w^2,\qquad
\gamma=2-4xy-x^2z,\quad u=1+xy,\quad w=\gamma u ,
\]
($q$ is \eqref{eq:normalization} for this $p$), and the map
\[
G\;=\;\Bigl(\gamma x,\;\;\frac{p(w)+2\gamma}{\gamma x},\;\;
\frac{q(w)+\gamma w}{(\gamma x)^2}\Bigr)^{\mathsf T}\colon\ \C^3\longrightarrow\C^3 .
\]

\begin{theorem}\label{thm:deg4}
$G$ is a polynomial map with components of degrees $4,11,12$; its Jacobian
determinant is identically $2$; and its generic fiber consists of exactly
$4$ points.
\end{theorem}

Here $b=-1$ and the components are in the displayed order, so
$\det J(G)=-2b=2$. The divisibilities are again explicit:
$p(w)+2\gamma=\gamma\bigl(6u-6\gamma u^2+\gamma^2u^3+2\bigr)$
and $q(w)+\gamma w=\gamma^2\bigl(\tfrac38\gamma^2u^4-2\gamma u^3+\tfrac32
u^2+u\bigr)$, and both bracketed factors vanish at $x=0$. The tangency
polynomial \eqref{eq:tangency} is the quartic
\[
\tfrac14 w^4-2w^3+3w^2-Xw+2Y ,
\]
whose leading coefficient is constant, so every target point admits four tangency
parameters counted with multiplicity.

Written out in full, $G=(g_1,g_2,g_3)^{\mathsf T}$ with
{\small
\begin{align*}
g_1 ={}& -x^3z-4x^2y+2x\\[2pt]
g_2 ={}& x^6y^3z^2+8x^5y^4z+3x^5y^2z^2+16x^4y^5+20x^4y^3z+3x^4yz^2+32x^3y^4+18x^3y^2z+x^3z^2\\
&+28x^2y^3+8x^2yz+16xy^2+2xz+2y\\[2pt]
g_3 ={}& \tfrac{3}{8}x^6y^4z^2+3x^5y^5z+\tfrac{3}{2}x^5y^3z^2+6x^4y^6+\tfrac{21}{2}x^4y^4z\\
&+\tfrac{9}{4}x^4y^2z^2+18x^3y^5+14x^3y^3z+\tfrac{3}{2}x^3yz^2+\tfrac{43}{2}x^2y^4\\
&+9x^2y^2z+\tfrac{3}{8}x^2z^2+14xy^3+3xyz+\tfrac{9}{2}y^2+\tfrac{1}{2}z
\end{align*}
}
All statements were verified exactly; the fiber count $4$ was confirmed at random
rational targets. The underlying curve $(p,q)^{\mathsf T}$ is a rational quartic with
\emph{two} cusps, at $w=2\pm\sqrt2$ (the common roots of $p'$ and $q'$), and, by
the genus count $\delta+\kappa=3$ for a rational quartic, one node; the Pl\"ucker
formula gives class $m=4\cdot3-2\cdot1-3\cdot2=4$, matching the fiber count. The
richer singularity structure of the curve produces a richer stratification than in
Theorem~\ref{thm:strat3}, which is worked out completely in
Appendix~\ref{app:G} (Theorem~\ref{thm:fibG}), and follows the pattern of
Proposition~\ref{prop:fibsizes} with $d=3$: double contact over smooth curve
points gives fibers of size $2$, the two cusps give size $1$, the node --- where
two branches each impose a double contact --- gives \emph{empty} fibers, and over
the hyperplane $v_1=0$, where the elimination quartic degenerates, the fiber has
$3$ points (or $1$ on a distinguished conic) --- the twist stratum invisible at
sweep level. Thus $G$ attains every fiber size in $\{4,3,2,1,0\}$.

\section{A general construction from parametrized hypersurfaces}\label{sec:dim4}

In this section we propose a general framework, valid in any dimension $n>2$:
one sweeps a tangent direction field on a parametrized hypersurface of dimension
$n-2$ in $\C^{n-1}$, pads by one multiplier variable, and twists by monomial
maps. The framework is set up in \S\ref{subsec:sweephyp} (where the branch
coefficients $L_k$ are listed explicitly for $n=4$ and $n=5$) and
\S\ref{subsec:forms} (which treats general $n$ only). We then specialize:
\S\ref{subsec:n4} is devoted to dimension four, with two counterexamples
$F_4$ and $F_5$, and \S\ref{subsec:n5} to dimension five, where the direction
field $\Delta=(1,w_1,w_1^2,w_2)^{\mathsf T}$ leads to a rigidity theorem for the middle
escape branch and, on the bottom branch, to a five-dimensional counterexample
$F_6$. Finally, \S\ref{subsec:curvefam} solves the curve-type direction field
$\Delta=(1,w_1,\dots,w_1^{n-2})^{\mathsf T}$ uniformly in every dimension --- the
normalization reduces to a $w_1$-family of constant-Jacobian maps of
$\C^{n-3}$ together with a free potential --- and produces a second
five-dimensional counterexample $F_7$, of geometric degree twelve.

\subsection{Sweeping a direction field on a hypersurface}\label{subsec:sweephyp}
The architecture sweeps a \emph{single} tangent direction field on a
parametrized hypersurface, and we set it up at once for arbitrary dimension
$n\ge3$. Let $w=(w_1,\dots,w_{n-2})$, let $X\in\C[w]^{\,n-1}$ be a column vector
of $n-1$ polynomials, viewed as a parametrized hypersurface
$X\colon\C^{n-2}\to\C^{n-1}$, and write $\partial_i=\partial/\partial w_i$, so
that $J(X)=(\partial_1X,\dots,\partial_{n-2}X)$ is an
$(n-1)\times(n-2)$ matrix whose column span is the tangent space of $X$ at $w$.
A \emph{polynomial tangent field} is a column vector $\Delta\in\C[w]^{\,n-1}$
with $\Delta(w)$ in the column span of $J(X)(w)$ for all $w$, equivalently the
\emph{tangency criterion}
\begin{equation}\label{eq:tangcrit}
\det\bigl(\Delta,\,J(X)\bigr)\;\equiv\;0 ,
\end{equation}
where $(\Delta,J(X))$ is the square matrix obtained by adjoining the column
$\Delta$. Under the standard identification
$\Lambda^{n-1}\C^{\,n-1}\cong\C$, the left-hand side is the exterior product
$\Delta\wedge\partial_1X\wedge\dots\wedge\partial_{n-2}X$ --- the
exterior-product form of the criterion referred to in the introduction. The padded sweep is
\[
F_0(x,\gamma,w)\;=\;\begin{pmatrix}\gamma x\\ X(w)+\gamma\Delta(w)\end{pmatrix}\colon
\ \C^n\to\C^n .
\]

\begin{lemma}\label{lem:JL}
Let $\Delta$ be a polynomial tangent field of $X$, i.e.\
$\det\bigl(\Delta,J(X)\bigr)\equiv0$. Then, with the variables ordered
$(x,\gamma,w_1,\dots,w_{n-2})$,
\[
\det J(F_0)\;=\;\gamma\,\det\bigl(\Delta,\;J(X)+\gamma J(\Delta)\bigr)
\;=\;\sum_{k=1}^{n-2}L_k\,\gamma^{\,k+1},
\]
where $L_k$ is the sum of the determinants $\det(\Delta,C_1,\dots,C_{n-2})$ in
which exactly $k$ of the columns $C_i$ are taken from $J(\Delta)$ and the rest
from $J(X)$.
\end{lemma}

We record the coefficients explicitly in the two dimensions used below. For
$n=4$ (parameters $w=(w_1,w_2)$, columns in $\C[w]^3$) we write $L:=L_1$ and
$M:=L_2$:
\[
\det J(F_0)=\gamma^2L+\gamma^3M,\qquad
L=\det(\Delta,\partial_1X,\partial_2\Delta)+\det(\Delta,\partial_1\Delta,\partial_2X),
\qquad
M=\det(\Delta,\partial_1\Delta,\partial_2\Delta).
\]
For $n=5$ (parameters $w=(w_1,w_2,w_3)$, columns in $\C[w]^4$),
$\det J(F_0)=\gamma^2L_1+\gamma^3L_2+\gamma^4L_3$ with
\begin{align*}
L_1&=\det(\Delta,\partial_1\Delta,\partial_2X,\partial_3X)
   +\det(\Delta,\partial_1X,\partial_2\Delta,\partial_3X)
   +\det(\Delta,\partial_1X,\partial_2X,\partial_3\Delta),\\
L_2&=\det(\Delta,\partial_1\Delta,\partial_2\Delta,\partial_3X)
   +\det(\Delta,\partial_1\Delta,\partial_2X,\partial_3\Delta)
   +\det(\Delta,\partial_1X,\partial_2\Delta,\partial_3\Delta),\\
L_3&=\det(\Delta,\partial_1\Delta,\partial_2\Delta,\partial_3\Delta).
\end{align*}
Note that $L_{n-2}$ involves $\Delta$ alone, and more generally the higher
coefficients are dominated by derivatives of $\Delta$: which of them can be made
a nonzero constant is decided by the direction field before the hypersurface is
chosen. All of these identities were verified by exact expansion.

\subsection{Prescribed direction fields and their tangency criterion}\label{subsec:forms}

In Lemma~\ref{lem:JL} the pair $(X,\Delta)$ is given together, with $\Delta$
tangent to $X$. In the constructions that follow we reverse the roles: we
\emph{prescribe} the direction field --- an explicit polynomial map
$\Delta\colon\C^{n-2}\to\C^{n-1}$, assigning to each parameter value $w$ the
direction that will be swept at the (yet unknown) point $X(w)$ --- and regard
the hypersurface $X$ as the unknown. For a fixed $\Delta$, the tangency
criterion \eqref{eq:tangcrit} then becomes a first-order partial differential
equation on the components of $X$, quasi-linear and multilinear in its first
derivatives; the normalization
adds the conditions on the coefficients $L_k$, in which $X$ appears with
progressively smaller weight as $k$ grows: $L_k$ is of degree $n-2-k$ in the
derivatives of $X$, and in particular the top coefficient $L_{n-2}$ involves
$\Delta$ alone. Which normalization branch is available is therefore decided by
the choice of $\Delta$ before any equation on $X$ is solved. This section works
with a general $n>2$ throughout; everything specific to $n=4$ and $n=5$ is
deferred to \S\ref{subsec:n4} and \S\ref{subsec:n5}.

Both the tangency criterion and Lemma~\ref{lem:JL} reduce to plain matrix
algebra, which we use to organize this analysis. Assume the column
$\Delta\in\C[w]^{\,n-1}$ is \emph{unimodular}, i.e.\ its entries generate the
unit ideal of $\C[w]$. Complete $\Delta$ to a matrix
$U\in\mathrm{SL}_{n-1}(\C[w])$ with first column $\Delta$, and let
\[
N\;\in\;\C[w]^{\,(n-2)\times(n-1)}
\]
be the matrix formed by the last $n-2$ rows of $U^{-1}$. Then $N\Delta=0$ (the
rows of $N$ are a free basis of the syzygies of $\Delta$), and for arbitrary
columns $V_1,\dots,V_{n-2}$,
\begin{equation}\label{eq:pairing}
\det\bigl(\Delta,V_1,\dots,V_{n-2}\bigr)
\;=\;\det\bigl(N\,(V_1,\dots,V_{n-2})\bigr),
\end{equation}
since $U^{-1}(\Delta,V_1,\dots,V_{n-2})$ has first column $e_1$ and
$\det U^{-1}=1$. Applying \eqref{eq:pairing} to the columns of
$J(X)+\gamma J(\Delta)$ and setting
\[
P\;:=\;N\,J(X),\qquad
\widetilde P\;:=\;N\,J(\Delta)
\qquad\bigl((n-2)\times(n-2)\ \text{matrices}\bigr),
\]
Lemma~\ref{lem:JL} becomes
\begin{equation}\label{eq:wedge}
\det J(F_0)\;=\;\gamma\,\det\bigl(P+\gamma\widetilde P\bigr),
\qquad\text{with}\qquad
\det P=0\ \ (\text{tangency}),
\end{equation}
and the branch coefficients $L_k$ of Lemma~\ref{lem:JL} are the coefficients of
$\det(P+\gamma\widetilde P)$ as a polynomial in $\gamma$: the tangency
$L_0=\det P\equiv0$, and, for $1\le k\le n-2$, $L_k$ is the sum of the
determinants obtained from $P$ by replacing exactly $k$ of its columns by the
corresponding columns of $\widetilde P$.

The normalization conditions now split into $n-2$ \emph{branches}, one for each
index $1\le k\le n-2$:
\[
\text{the \emph{$L_k$-branch}:}\qquad
L_j\equiv0\ (j\neq k),\qquad L_k\equiv\text{const}\neq0,
\qquad\text{giving}\qquad \det J(F_0)=L_k\,\gamma^{\,k+1} .
\]
(For $n=4$, where $L=L_1$ and $M=L_2$ as in \S\ref{subsec:sweephyp}, the two
branches are called the \emph{$L$-branch} and the \emph{$M$-branch}.)
Which branches can be realized is constrained by $\widetilde P=N\,J(\Delta)$,
which depends on $\Delta$ alone: since every term of $L_k$ contains $k$ columns
of $\widetilde P$, one has $L_j\equiv0$ automatically for all
$j>\operatorname{rank}\widetilde P$, so the rank of $\widetilde P$ bounds the
highest available branch. At the two extremes, a direction field depending on a
single parameter forces $\operatorname{rank}\widetilde P\le1$ (only the bottom
branch survives), while the tautological field with $\widetilde P$ the identity
opens the top branch $k=n-2$. Mixed direction fields, whose $\widetilde P$ has
intermediate rank, first occur for $n=5$ and are the subject of
\S\ref{subsec:n5}.

Two structural remarks apply for every $n$. First, on the locus where the first
row of $P$ is nonzero, tangency ($\det P\equiv0$, i.e.\ $P$ singular) expresses
the remaining rows through rational multipliers, and the equality of mixed
partial derivatives of the entries of $X$ converts the row relations defining
$P$ into \emph{scalar potentials}; the branch equations then become a system of
first-order partial differential equations on the potentials. Second, when
$\Delta$ has a constant entry, the potentials can be chosen globally. Write
$\Delta=(\delta_1,\dots,\delta_{n-1})^{\mathsf T}$ with $\delta_1=1$, take the
syzygy matrix $N$ with rows $-\,\delta_{i+1}e_1^{\mathsf T}+e_{i+1}^{\mathsf T}$
($1\le i\le n-2$), and define the \emph{potentials}
\[
G_i\;:=\;X_i-\delta_i\,X_1\qquad(i=2,\dots,n-1),
\qquad
\Gamma\;:=\;(G_2,\dots,G_{n-1})^{\mathsf T},
\qquad
\widehat\Delta\;:=\;(\delta_2,\dots,\delta_{n-1})^{\mathsf T}.
\]
Differentiating by the product rule,
$\partial_jG_i=\partial_jX_i-\delta_i\,\partial_jX_1-X_1\,\partial_j\delta_i$,
which in matrix form is
\[
P\;=\;N\,J(X)\;=\;J(\Gamma)+X_1\,J(\widehat\Delta),
\qquad\qquad
\widetilde P\;=\;N\,J(\Delta)\;=\;J(\widehat\Delta)
\]
(for the second identity, the row $-\delta_{i+1}e_1^{\mathsf T}+e_{i+1}^{\mathsf T}$
applied to $\partial_j\Delta$ gives
$\partial_j\delta_{i+1}-\delta_{i+1}\partial_j\delta_1=\partial_j\delta_{i+1}$,
since $\delta_1=1$ is constant). The whole branch analysis therefore takes
place in the $(n-2)\times(n-2)$ pencil
\[
P+\gamma\widetilde P\;=\;J(\Gamma)+\mu\,J(\widehat\Delta),
\qquad
\mu\;=\;X_1+\gamma ,
\]
in which the hypersurface enters only through the potentials $\Gamma$ and the
scalar $X_1$; this device is used systematically in \S\ref{subsec:n5}
and \S\ref{subsec:curvefam}.

\subsection{The twist in dimension $n$: the stage equations}\label{subsec:stage}

All the higher-dimensional maps below are produced by one and the same
procedure, which converts a normalized sweep into a Keller map. Since we use
it repeatedly, we record it once, together with the equations that have to be
solved. The input is a padded sweep $F_0=(\gamma x,\,S)^{\mathsf T}$,
$S=X+\gamma\Delta$, on a branch of \S\ref{subsec:forms}:
$\det J(F_0)=c\,\gamma^{\,k+1}$. The output has the shape
\[
F\;=\;\Bigl(C,\;\frac{S_1}{C^{\,e_1}},\;\dots,\;
\frac{S_{n-1}}{C^{\,e_{n-1}}}\Bigr)^{\mathsf T},
\qquad C=\gamma x ,
\]
after $\gamma$ and $w$ are expressed through the source variables by a
monomial map and a stage. Three layers of choices are involved.

\emph{(i) Discrete data.} Positive integers: weights $d_j$ for the parameters
($\operatorname{wt}w_j=d_j$), twist exponents $e_i$ for the target components,
and $x$-degrees $m_j$ for the source monomials. They must satisfy
\begin{equation}\label{eq:discrete}
\operatorname{wt}(X_i)\ge e_i\ \ (\text{every monomial}),\qquad
1+\operatorname{wt}(\Delta_i)\ge e_i,\qquad
\sum_j m_j=\sum_i e_i,\qquad
\sum_j d_j+(k+1)=\sum_i e_i .
\end{equation}
The first two conditions make
$E_i:=S_i\bigl(\gamma,\;\gamma^{d_1}u_1,\dots,\gamma^{d_{n-2}}u_{n-2}\bigr)
/\gamma^{\,e_i}$ a \emph{polynomial} in $(\gamma,u)$; the last two make the
$x$- and $\gamma$-powers cancel in the determinant chain below.

\emph{(ii) The stage.} On the source $\C^n\ni(x,y_1,\dots,y_{n-1})$ put
$v_j=x^{m_j}y_j$ and let the stage be
\[
(\gamma,u_1,\dots,u_{n-2})^{\mathsf T}
\;=\;\theta^0+L\,v+\textstyle\sum_\mu q_\mu\,\mu(v):
\]
a base point $\theta^0$, an invertible linear part $L$ with columns
$\mathrm{col}_1,\dots,\mathrm{col}_{n-1}$, and finitely many correction
vectors $q_\mu\in\C^{n-1}$ attached to monomials $\mu(v)$ of degree $\ge2$.
Let $J\subseteq\{1,\dots,n-1\}$ be the set of indices of the variables $v_j$
occurring in some correction monomial.

\begin{lemma}\label{lem:stage}
If every correction vector lies in
$V:=\mathrm{span}\{\mathrm{col}_j:j\notin J\}$, then the stage is a polynomial
automorphism of $\C^{n-1}$ with constant Jacobian determinant $\det L$.
\end{lemma}

\begin{proof}
The Jacobian is $L+\sum_\mu q_\mu\,\nabla\mu(v)$; in the multilinear expansion
of its determinant, any term replaces a nonempty set of columns
$\mathrm{col}_j$ with $j\in J$ (since $\partial_j\mu\neq0$ forces $j\in J$) by
vectors in $V$; together with the untouched columns
$\{\mathrm{col}_j:j\notin J\}$ this places more than $\dim V$ columns inside
$V$, so every such term vanishes. For the inverse: $L^{-1}q_\mu\in
\mathrm{span}\{e_j:j\notin J\}$, so the $J$-coordinates of $v$ are affine in
the stage values; substituting them into the corrections (which involve only
those coordinates) recovers the remaining coordinates polynomially.
\end{proof}

\emph{(iii) The side conditions.} $F$ is a polynomial map iff
$x^{\,e_i}\mid E_i\circ\text{stage}$ for every $i$; since $v_j$ has
$x$-degree $m_j$, this is the finite system
\begin{equation}\label{eq:side}
\bigl[\mu\bigr]\bigl(E_i\circ\text{stage}\bigr)\;=\;0
\qquad\text{for every monomial $\mu(v)$ of $x$-degree}<e_i,\ i=1,\dots,n-1 .
\end{equation}
Ordered by the $x$-degree $r$ of $\mu$, the system is \emph{triangular}:
\begin{itemize}
\item $r=0$: $E_i(\theta^0)=0$ for all $i$ --- equivalently $S(\gamma_0,w^0)=0$
at the base parameters $w_j^{\,0}=\gamma_0^{\,d_j}u_j^0$: the swept
hypersurface passes through the origin. These $n-1$ equations are absorbed
into the free data of $X$.
\item $r=1$: $\nabla E_i(\theta^0)\cdot\mathrm{col}_j=0$ for each $v_j$ of
$x$-degree $1$ and each $i$ with $e_i\ge2$: the low columns lie in common
kernels of gradients.
\item $r\ge2$: for each column of $x$-degree $r$, kernel conditions
$\nabla E_i\cdot\mathrm{col}_j=0$ ($e_i>r$); and for each correction monomial
$\mu$ of $x$-degree $r$, the inhomogeneous linear equation
\[
\nabla E_i\cdot q_\mu\;=\;-\bigl[\mu\bigr]
\bigl(E_i\circ(\text{affine part of the stage})\bigr)
\qquad(e_i>r),
\]
whose right-hand side is a Taylor coefficient of $E_i$ along the columns
already chosen.
\end{itemize}
Every unknown --- the base point, a column, a correction vector --- enters
\eqref{eq:side} for the first time \emph{linearly}, so the hierarchy is solved
by finite-dimensional linear algebra, order by order; the only nondegeneracy
required is the linear independence of the gradients involved at each step,
together with the invertibility of the assembled $L$ and the span condition of
Lemma~\ref{lem:stage}. The Jacobian determinants then multiply to
\begin{equation}\label{eq:chain}
\det J(F)\;=\;
\underbrace{x^{\sum m_j}}_{\text{monomials}}\cdot
\underbrace{\det L}_{\text{stage}}\cdot
\underbrace{\gamma^{\sum d_j}}_{\text{scalings}}\cdot
\underbrace{c\,\gamma^{\,k+1}}_{\det J(F_0)}\cdot
\underbrace{(\gamma x)^{-\sum e_i}}_{\text{twist}}
\;=\;c\,\det L ,
\end{equation}
a nonzero constant. The discrete data of the six maps of this paper are:
\begin{center}
\begin{tabular}{lcccccc}
\hline
map & $n$ & $\det J(F_0)$ & $(d_j)$ & $(e_i)$ & $(m_j)$ & $\det J(F)=\pm c\det L$\\
\hline
$F$   & $3$ & $2\gamma^2$ & $(1)$     & $(1,2)$     & $(1,2)$     & $-2$\\
$G$   & $3$ & $2\gamma^2$ & $(1)$     & $(1,2)$     & $(1,2)$     & $2$\\
$F_4$ & $4$ & $2\gamma^2$ & $(1,1)$   & $(1,2,1)$   & $(1,1,2)$   & $-\tfrac{44}{9}$\\
$F_5$ & $4$ & $\gamma^3$  & $(1,2)$   & $(1,2,3)$   & $(1,2,3)$   & $\tfrac{160}{29}$\\
$F_6$ & $5$ & $\gamma^2$  & $(1,3,4)$ & $(1,2,3,4)$ & $(1,2,3,4)$ & $-290$\\
$F_7$ & $5$ & $\gamma^2$  & $(1,3,4)$ & $(1,2,3,4)$ & $(1,2,3,4)$ & $119377$\\
\hline
\end{tabular}
\end{center}
(the sign for $F$ comes from the traditional component reversal). All the
stage displays in the following subsections are solutions of
\eqref{eq:side} found exactly this way.

\subsection{The case $n=4$}\label{subsec:n4}

\subsubsection{Specialization I: $\Delta=(1,w_1,w_1^2)^{\mathsf T}$ and the map $F_4$}
\label{subsec:specI}

The column $\Delta=(1,w_1,w_1^2)^{\mathsf T}$ is unimodular, with the natural
syzygy matrix of \S\ref{subsec:forms} (rows
$-\delta_{i+1}e_1^{\mathsf T}+e_{i+1}^{\mathsf T}$):
\[
N=\begin{pmatrix}-w_1&1&0\\[1pt] -w_1^2&0&1\end{pmatrix},
\qquad N\Delta=0,\qquad
\det(\Delta,V,W)=\det\bigl(N\,(V,W)\bigr)\ \ \forall V,W .
\]
Since $\Delta$ depends only on $w_1$, its independent partial derivatives are
$\Delta'=\partial_1\Delta=(0,1,2w_1)^{\mathsf T}$ and
$\Delta''=\partial_1^2\Delta=(0,0,2)^{\mathsf T}$ (all higher ones vanish), and
$(\Delta,\Delta',\Delta'')$ is a frame of constant volume
$\det(\Delta,\Delta',\Delta'')=2$; against the syzygy matrix,
\[
N\Delta=0,\qquad
N\Delta'=\begin{pmatrix}1\\ 2w_1\end{pmatrix},\qquad
N\Delta''=\begin{pmatrix}0\\ 2\end{pmatrix}.
\]
For a ruled surface $X=A(w_1)+w_2B(w_1)$ --- so that
$\partial_1X=A'(w_1)+w_2B'(w_1)$ and $\partial_2X=B(w_1)$ --- expand the data
in the frame,
\[
B=b_1\Delta+b_2\Delta'+b_3\Delta'',\qquad
A'=a_1\Delta+a_2\Delta'+a_3\Delta''\qquad(a_i,b_i\in\C[w_1]),
\]
so that, using $\Delta'''=0$,
$B'=b_1'\Delta+(b_1+b_2')\Delta'+(b_2+b_3')\Delta''$. The product
$P=N\,J(X)=(N\,\partial_1X,\;N\,\partial_2X)$ is now explicit: abbreviating
$c_2=a_2+w_2(b_1+b_2')$ and $c_3=a_3+w_2(b_2+b_3')$,
\[
P=\begin{pmatrix}c_2 & b_2\\ 2(w_1c_2+c_3) & 2(w_1b_2+b_3)\end{pmatrix},
\qquad
\widetilde P=N\,J(\Delta)=\begin{pmatrix}1&0\\ 2w_1&0\end{pmatrix},
\]
and therefore
\[
\det\bigl(P+\gamma\widetilde P\bigr)\;=\;2\,(c_2b_3-c_3b_2)\;+\;2\,\gamma\,b_3 .
\]
The coefficient $M$ of $\gamma^2$ vanishes \emph{automatically} --- this
$\Delta$ lives on the $L$-branch --- the normalization $L\equiv2$ forces
$b_3\equiv1$, and the tangency, i.e.\ the identical vanishing in $w_2$ of the
$\gamma$-free term, reduces to
\[
b_3=1,\qquad b_1=b_2^2-b_2',\qquad a_2=a_3b_2 .
\]
Thus, on the ruled ansatz, the branch data are \emph{three free polynomials}
$b_2,a_3,a_1\in\C[w_1]$ (together with the constant of integration of $A$):
the relations displayed determine $b_3,b_1,a_2$, the curve $A$ is recovered
from $A'$ by integration --- which stays polynomial --- and \emph{every}
such choice satisfies $\det J(F_0)=2\gamma^2$ (verified by exact expansion on
randomized data). The finitely many side conditions of the twist then select
special members of this family. (In the language of \S\ref{subsec:curvefam},
this is the case $n=4$ of the curve-type family with $G_2$ linear in $w_2$:
the three free univariate polynomials correspond to the coefficients of
$G_2=\alpha(w_1)w_2+\beta(w_1)$ and to $\varphi(w_1)$ in
$H_3=cw_2+\varphi$.)
Crucially there are solutions with $b_2\not\equiv0$: the surface is then ruled but
\emph{not} a cylinder, and the second parameter $w_2$ enters every swept
coordinate. (Cylindrical solutions $b_1=b_2=0$ produce only shears of
three-dimensional examples.)

Take $b_2=w_1$, $a_3=1$ (so $a_2=w_1$, $b_1=w_1^2-1$) and
$a_1=\tfrac13-\tfrac83w_1$. Integrating, the curve and ruling are
\[
A(w_1)=\Bigl(\tfrac13w_1-\tfrac43w_1^2,\;\;
\tfrac23w_1^2-\tfrac89w_1^3,\;\;
2w_1+\tfrac79w_1^3-\tfrac23w_1^4\Bigr)^{\mathsf T},\qquad
B(w_1)=\bigl(w_1^2-1,\;w_1^3,\;w_1^4+w_1^2+2\bigr)^{\mathsf T},
\]
so the surface $X=A+w_2B=(p,q,r)^{\mathsf T}$ and the tangent field $\Delta$ are, explicitly,
\begin{align*}
p(w_1,w_2)&=\tfrac13w_1-\tfrac43w_1^2+w_2\,(w_1^2-1),\\
q(w_1,w_2)&=\tfrac23w_1^2-\tfrac89w_1^3+w_2\,w_1^3,\\
r(w_1,w_2)&=2w_1+\tfrac79w_1^3-\tfrac23w_1^4+w_2\,(w_1^4+w_1^2+2),\\
\Delta(w_1,w_2)&=(1,\;w_1,\;w_1^2)^{\mathsf T}
\qquad(\text{independent of }w_2).
\end{align*}
The tangency identity $\det(\Delta,J(X))\equiv0$ and the
normalization $L\equiv2$, $M\equiv0$ hold for this pair by construction (and were
verified by expansion).
Define on $\C^4\ni(x,y,z,t)$:
\[
\gamma=1+xy+x^2t,\quad u_1=1+xz,\quad
u_2=-\tfrac79+\tfrac{22}9xy+\tfrac83xz,\quad
w_1=\gamma u_1,\quad w_2=\gamma u_2,\quad C=\gamma x,
\]
\[
S_i=A_i(w_1)+w_2B_i(w_1)+\gamma\,\Delta_i(w_1)\quad(i=1,2,3),
\]
\[
\boxed{\;F_4\;=\;\Bigl(C,\;\;\frac{S_1}{C},\;\;\frac{S_2}{C^2},\;\;\frac{S_3}{C}\Bigr)^{\mathsf T}.\;}
\]
In the format of the composition diagram of \S\ref{subsec:twist}, $F_4$ is the
composite
\[
\begin{pmatrix}x\\ y\\ z\\ t\end{pmatrix}
\xrightarrow{\ \text{monomial}\ }
\begin{pmatrix}x\\ xy\\ xz\\ x^2t\end{pmatrix}
\xrightarrow{\ \text{affine}\ }
\begin{pmatrix}x\\ \gamma\\ u_1\\ u_2\end{pmatrix}
\xrightarrow{\ \substack{w_1=\gamma u_1\\[1pt] w_2=\gamma u_2}\ }
\begin{pmatrix}x\\ \gamma\\ w_1\\ w_2\end{pmatrix}
\xrightarrow{\ \text{padded sweep}\ }
\begin{pmatrix}\gamma x\\ S_1\\ S_2\\ S_3\end{pmatrix}
\xrightarrow{\ \text{twist}\ }
\begin{pmatrix}C\\ S_1/C\\ S_2/C^2\\ S_3/C\end{pmatrix},
\]
whose Jacobian determinants multiply to the constant computed after
Theorem~\ref{thm:F4} below.

\begin{theorem}\label{thm:F4}
$F_4$ is a polynomial map $\C^4\to\C^4$ with components of degrees $4,11,12,21$;
its Jacobian determinant is identically $-\tfrac{44}{9}$; and its generic fiber
consists of exactly $5$ points.
\end{theorem}

Consequently $F_4$ is a non-injective Keller map
in dimension four which is not equivalent, under composition with polynomial
automorphisms on either side, to $\Phi\times\mathrm{id}$ for any
three-dimensional Keller map $\Phi$ of geometric degree $\neq5$: geometric degree
is invariant under such compositions, and $\deg(\Phi\times\mathrm{id})
=\deg\Phi$. In particular $F_4$ is not a trivial extension of Alp\"oge's example
or of the map $G$. Whether $F_4$ is equivalent to $\Phi\times\mathrm{id}$ for a
degree-\emph{five} member $\Phi$ of Gallagher's family \cite{Gallagher2026} is
left open; the fiber stratification of a product has product form, so the
stratification deferred to Appendix~\ref{app:F4} is expected to decide this.

\begin{proof}[Proof sketch]
The degrees and divisibilities are read off the construction, and
$\det J(F_4)\equiv-\tfrac{44}9$ is the chain factorization through the stage
(Lemma~\ref{lem:stage}); the constancy was also confirmed at random rational
points. For the fibers we use the elimination scheme that recurs for $F_5$,
$F_6$ and $F_7$: eliminate $\gamma$ through the first sweep equation, then
eliminate the remaining transverse parameter by a resultant, leaving a
univariate polynomial in $w_1$ certified by exact gcd computations. Here, let
$v_1\neq0$, set $(Y_1,Y_2,Y_3)=(v_1v_2,\,v_1^2v_3,\,v_1v_4)$, and substitute
$\gamma=Y_1-A_1-w_2B_1$ (the first sweep equation; $\Delta_1=1$) into the
other two. The resulting equations are \emph{linear} in $w_2$, with leading
coefficients $w_1$ and $2w_1^2+2$ having no common zero, and their resultant
with respect to $w_2$ is, up to sign, the compatibility determinant
\[
R(w_1;Y)\;=\;\det\bigl(A(w_1)-Y,\;B(w_1),\;\Delta(w_1)\bigr)
\;=\;\tfrac29w_1^5+\tfrac29w_1^4+\bigl(\tfrac89+Y_1\bigr)w_1^3
-\bigl(\tfrac43+2Y_2\bigr)w_1^2+(2Y_1+Y_3)\,w_1-2Y_2 ,
\]
a quintic with \emph{constant} leading coefficient $\tfrac29$ --- the tangency
quintic of the swept surface, playing the role of the tangency polynomial
\eqref{eq:tangency}. Each root carries a unique $(w_2,\gamma)$ (the two
leading coefficients are coprime), and each root with $\gamma\neq0$ lifts to
exactly one fiber point via $x=v_1/\gamma$ and the triangular stage, which
recovers $(y,z,t)$ rationally. At the witness target
$v=(\tfrac23,-\tfrac15,\tfrac12,3)$, exact univariate gcd certificates show
that $R$ is squarefree and coprime to the $w_2$-resultant of the second sweep
equation with $\gamma=0$, so that fiber has exactly $5$ points, and by
openness of these conditions the generic count is $5$.
\end{proof}

Written out in full, $F_4=(F_{4,1},F_{4,2},F_{4,3},F_{4,4})^{\mathsf T}$ with
{\small
\begin{align*}
F_{4,1} ={}& x^3t+x^2y+x\\[2pt]
F_{4,2} ={}&\tfrac{22}{9}x^6yz^2t^2+\tfrac{8}{3}x^6z^3t^2+\tfrac{44}{9}x^5y^2z^2t+\tfrac{16}{3}x^5yz^3t+\tfrac{44}{9}x^5yzt^2+\tfrac{41}{9}x^5z^2t^2+\tfrac{22}{9}x^4y^3z^2\\
&+\tfrac{8}{3}x^4y^2z^3+\tfrac{88}{9}x^4y^2zt+14x^4yz^2t+\tfrac{16}{3}x^4z^3t+\tfrac{22}{9}x^4yt^2+\tfrac{10}{9}x^4zt^2+\tfrac{44}{9}x^3y^3z+\tfrac{85}{9}x^3y^2z^2\\
&+\tfrac{16}{3}x^3yz^3+\tfrac{44}{9}x^3y^2t+12x^3yzt+\tfrac{70}{9}x^3z^2t-\tfrac{7}{9}x^3t^2+\tfrac{22}{9}x^2y^3+\tfrac{98}{9}x^2y^2z+\tfrac{92}{9}x^2yz^2+\tfrac{8}{3}x^2z^3\\
&+\tfrac{10}{3}x^2yt-\tfrac{4}{9}x^2zt+\tfrac{37}{9}xy^2+\tfrac{40}{9}xyz+\tfrac{29}{9}xz^2-\tfrac{26}{9}xt-\tfrac{26}{9}y-\tfrac{35}{9}z\\[2pt]
F_{4,3} ={}& \tfrac{22}{9}x^6yz^3t^2+\tfrac{8}{3}x^6z^4t^2+\tfrac{44}{9}x^5y^2z^3t+\tfrac{16}{3}x^5yz^4t+\tfrac{22}{3}x^5yz^2t^2+\tfrac{65}{9}x^5z^3t^2+\tfrac{22}{9}x^4y^3z^3\\
&+\tfrac{8}{3}x^4y^2z^4+\tfrac{44}{3}x^4y^2z^2t+\tfrac{58}{3}x^4yz^3t+\tfrac{16}{3}x^4z^4t+\tfrac{22}{3}x^4yzt^2+\tfrac{17}{3}x^4z^2t^2+\tfrac{22}{3}x^3y^3z^2\\
&+\tfrac{109}{9}x^3y^2z^3+\tfrac{16}{3}x^3yz^4+\tfrac{44}{3}x^3y^2zt+26x^3yz^2t+\tfrac{122}{9}x^3z^3t+\tfrac{22}{9}x^3yt^2+\tfrac{1}{3}x^3zt^2+\tfrac{22}{3}x^2y^3z\\
&+\tfrac{61}{3}x^2y^2z^2+16x^2yz^3+\tfrac{8}{3}x^2z^4+\tfrac{44}{9}x^2y^2t+\tfrac{46}{3}x^2yzt+\tfrac{26}{3}x^2z^2t-\tfrac{7}{9}x^2t^2+\tfrac{22}{9}xy^3+15xy^2z\\
&+16xyz^2+\tfrac{19}{3}xz^3+\tfrac{10}{3}xyt-2xzt+\tfrac{37}{9}y^2+\tfrac{16}{3}yz+\tfrac{11}{3}z^2-\tfrac{22}{9}t\\[2pt]
F_{4,4} ={}& \tfrac{22}{9}x^{12}yz^4t^4+\tfrac{8}{3}x^{12}z^5t^4+\tfrac{88}{9}x^{11}y^2z^4t^3+\tfrac{32}{3}x^{11}yz^5t^3+\tfrac{88}{9}x^{11}yz^3t^4+\tfrac{89}{9}x^{11}z^4t^4\\
&+\tfrac{44}{3}x^{10}y^3z^4t^2+16x^{10}y^2z^5t^2+\tfrac{352}{9}x^{10}y^2z^3t^3+\tfrac{148}{3}x^{10}yz^4t^3+\tfrac{32}{3}x^{10}z^5t^3+\tfrac{88}{9}x^9y^4z^4t\\
&+\tfrac{32}{3}x^9y^3z^5t+\tfrac{44}{3}x^{10}yz^2t^4+\tfrac{116}{9}x^{10}z^3t^4+\tfrac{176}{3}x^9y^3z^3t^2+\tfrac{266}{3}x^9y^2z^4t^2+32x^9yz^5t^2\\
&+\tfrac{22}{9}x^8y^5z^4+\tfrac{8}{3}x^8y^4z^5+\tfrac{176}{3}x^9y^2z^2t^3+\tfrac{272}{3}x^9yz^3t^3+\tfrac{350}{9}x^9z^4t^3+\tfrac{352}{9}x^8y^4z^3t+\tfrac{620}{9}x^8y^3z^4t\\
&+32x^8y^2z^5t+\tfrac{88}{9}x^9yzt^4+6x^9z^2t^4+88x^8y^3z^2t^2+\tfrac{584}{3}x^8y^2z^3t^2+\tfrac{394}{3}x^8yz^4t^2+16x^8z^5t^2\\
&+\tfrac{88}{9}x^7y^5z^3+\tfrac{59}{3}x^7y^4z^4+\tfrac{32}{3}x^7y^3z^5+\tfrac{352}{9}x^8y^2zt^3+\tfrac{248}{3}x^8yz^2t^3+\tfrac{440}{9}x^8z^3t^3+\tfrac{176}{3}x^7y^4z^2t\\
&+\tfrac{1520}{9}x^7y^3z^3t+146x^7y^2z^4t+32x^7yz^5t+\tfrac{22}{9}x^8yt^4-\tfrac{4}{9}x^8zt^4+\tfrac{176}{3}x^7y^3zt^2+212x^7y^2z^2t^2\\
&+\tfrac{616}{3}x^7yz^3t^2+\tfrac{172}{3}x^7z^4t^2+\tfrac{44}{3}x^6y^5z^2+52x^6y^4z^3+\tfrac{482}{9}x^6y^3z^4+16x^6y^2z^5+\tfrac{88}{9}x^7y^2t^3\\
&+\tfrac{112}{3}x^7yzt^3+20x^7z^2t^3+\tfrac{352}{9}x^6y^4zt+200x^6y^3z^2t+264x^6y^2z^3t+\tfrac{1120}{9}x^6yz^4t+\tfrac{32}{3}x^6z^5t\\
&-\tfrac{7}{9}x^7t^4+\tfrac{44}{3}x^6y^3t^2+\tfrac{344}{3}x^6y^2zt^2+\tfrac{1354}{9}x^6yz^2t^2+\tfrac{655}{9}x^6z^3t^2+\tfrac{88}{9}x^5y^5z+\tfrac{194}{3}x^5y^4z^2\\
&+\tfrac{968}{9}x^5y^3z^3+\tfrac{604}{9}x^5y^2z^4+\tfrac{32}{3}x^5yz^5+\tfrac{20}{3}x^6yt^3-\tfrac{40}{9}x^6zt^3+\tfrac{88}{9}x^5y^4t+\tfrac{1040}{9}x^5y^3zt\\
&+\tfrac{2168}{9}x^5y^2z^2t+\tfrac{554}{3}x^5yz^3t+\tfrac{338}{9}x^5z^4t+\tfrac{74}{3}x^5y^2t^2+\tfrac{452}{9}x^5yzt^2+\tfrac{287}{9}x^5z^2t^2+\tfrac{22}{9}x^4y^5\\
&+\tfrac{116}{3}x^4y^4z+\tfrac{994}{9}x^4y^3z^2+\tfrac{1007}{9}x^4y^2z^3+40x^4yz^4+\tfrac{8}{3}x^4z^5-\tfrac{34}{9}x^5t^3+\tfrac{236}{9}x^4y^3t+\tfrac{1024}{9}x^4y^2zt\\
&+\tfrac{382}{3}x^4yz^2t+\tfrac{454}{9}x^4z^3t+\tfrac{52}{9}x^4yt^2-\tfrac{47}{9}x^4zt^2+9x^3y^4+\tfrac{532}{9}x^3y^3z+\tfrac{859}{9}x^3y^2z^2+\tfrac{542}{9}x^3yz^3\\
&+\tfrac{83}{9}x^3z^4+\tfrac{206}{9}x^3y^2t+\tfrac{346}{9}x^3yzt+\tfrac{250}{9}x^3z^2t-\tfrac{17}{3}x^3t^2+\tfrac{40}{3}x^2y^3+\tfrac{131}{3}x^2y^2z+\tfrac{404}{9}x^2yz^2\\
&+\tfrac{41}{3}x^2z^3+\tfrac{10}{3}x^2yt+\tfrac{10}{9}x^2zt+9xy^2+\tfrac{142}{9}xyz+\tfrac{89}{9}xz^2-\tfrac{28}{9}xt+\tfrac{20}{3}y+\tfrac{29}{3}z
\end{align*}
}

The proof of the Jacobian statement is the verified chain factorization
$x^4\cdot(-\tfrac{22}9)\cdot\gamma^2\cdot2\gamma^2\cdot(\gamma x)^{-4}
=-\tfrac{44}9$, in which the factor $2\gamma^2$ is Lemma~\ref{lem:JL} for the
surface above ($L=2$, $M=0$; a polynomial identity), the factor $\gamma^2$ comes
from the scaling $(w_1,w_2)=(\gamma u_1,\gamma u_2)$, $x^4$ from the monomial
substitution $(v_1,v_2,v_3)=(xy,xz,x^2t)$ with $x$-degrees $(1,1,2)$, and
$-\tfrac{22}9$ is the determinant of the affine stage. The five side conditions
(divisibility of $S_1,S_3$ by $C$ and of $S_2$ by $C^2$) determine the constants
$(\gamma_0,u_1^0,u_2^0,\alpha_0,\alpha_1)=(1,1,-\tfrac79,\tfrac13,-\tfrac83)$;
the grading $(1,1,2)$ leaves the affine stage invertible. The fiber count was
established by exact
elimination: the tangency equation of the swept ruled family has degree $7$,
always divisible by the spurious factor $w_1^2-1$, leaving a squarefree quintic;
all five roots lift, as confirmed by explicit fibers with residuals below
$10^{-9}$, e.g. the target of $(0.7,-0.4,0.3,0.5)$ has five preimages, three real
and one conjugate pair.

\subsubsection{Specialization II: $\Delta=(1,w_1,w_2)^{\mathsf T}$ and a map of geometric degree ten}
\label{subsec:specII}

The column $\Delta=(1,w_1,w_2)^{\mathsf T}$ is unimodular with the natural
syzygy matrix
\[
N=\begin{pmatrix}-w_1&1&0\\ -w_2&0&1\end{pmatrix},
\qquad N\Delta=0 ,
\]
and now the partial derivatives of $\Delta$ are the \emph{constant} columns
$\partial_1\Delta=e_2$ and $\partial_2\Delta=e_3$: the triple
$(\Delta,\partial_1\Delta,\partial_2\Delta)$ is a frame of constant volume
$\det(\Delta,e_2,e_3)=1$, with $N\Delta=0$, $Ne_2=(1,0)^{\mathsf T}$,
$Ne_3=(0,1)^{\mathsf T}$, so that $\widetilde P=N\,J(\Delta)=I_2$, the identity
matrix. Exactly as in \S\ref{subsec:specI}, expand the columns of $J(X)$ in the
frame,
\[
\partial_1X=a_1\Delta+a_2e_2+a_3e_3,\qquad
\partial_2X=b_1\Delta+b_2e_2+b_3e_3
\qquad(a_i,b_i\in\C[w_1,w_2]);
\]
comparing first entries forces $a_1=\partial_1X_1$ and $b_1=\partial_2X_1$, so
the first coefficients are automatically integrable, with $X_1$ itself as their
potential. The products with $N$ are again explicit:
\[
P=N\,J(X)=\begin{pmatrix}a_2 & b_2\\ a_3 & b_3\end{pmatrix},
\qquad
\det\bigl(P+\gamma\widetilde P\bigr)
\;=\;(a_2b_3-a_3b_2)\;+\;\gamma\,(a_2+b_3)\;+\;\gamma^2 .
\]
Thus $M\equiv1$ identically and $L=a_2+b_3$, the trace of $P$: this $\Delta$
lives on the $M$-branch, with normalization $L\equiv0$ and
$\det J(F_0)=\gamma^3$.

The branch is solved completely by combining the two algebraic conditions with
the integrability of the expansion. The tangency and the normalization
$L\equiv0$ give two equations,
\[
a_2b_3-a_3b_2\;=\;0
\qquad\text{and}\qquad
a_2+b_3\;=\;0 ,
\]
which say that $P$ is trace-free of rank $\le1$; this is parametrized by two
scalars $\rho,\lambda$:
\[
(a_2,\;b_2,\;a_3,\;b_3)\;=\;(-\rho\lambda,\;\lambda,\;-\rho^2\lambda,\;\rho\lambda).
\]
The equality of mixed partials $\partial_2(\partial_1X)=\partial_1(\partial_2X)$,
expanded in the frame (using $\partial_2\Delta=e_3$ and $\partial_1\Delta=e_2$,
the $e_i$ being constant), gives three scalar equations:
\[
\partial_2a_1=\partial_1b_1\ \ (\text{automatic}),\qquad
\partial_2a_2=b_1+\partial_1b_2,\qquad
a_1+\partial_2a_3=\partial_1b_3 .
\]
The last two express the gradient of $X_1$ through $(\rho,\lambda)$:
\[
\partial_2X_1=\partial_2a_2-\partial_1b_2=-\,\partial_2(\rho\lambda)-\partial_1\lambda,
\qquad
\partial_1X_1=\partial_1b_3-\partial_2a_3=\partial_1(\rho\lambda)+\partial_2(\rho^2\lambda),
\]
and the compatibility $\partial_1\partial_2X_1=\partial_2\partial_1X_1$ of this
gradient is the single \emph{linear} wave-type equation
\begin{equation}\label{eq:wave}
\lambda_{11}+2(\rho\lambda)_{12}+(\rho^2\lambda)_{22}\;=\;0 ,
\end{equation}
after which $X_1$ is recovered by integration,
and then the potentials
$G_2=X_2-w_1X_1$ and $G_3=X_3-w_2X_1$ --- whose gradients are, by
\S\ref{subsec:forms}, the rows of $P-X_1\widetilde P$ --- by two further
integrations.

Solutions of \eqref{eq:wave} are produced in two steps. First fix the
multiplier $\rho\in\C[w_1,w_2]$; since the rank-one parametrization above is
linear in $\lambda$ for fixed $\rho$, equation \eqref{eq:wave} is then
\emph{linear} in $\lambda$, with polynomial coefficients (a degenerate
second-order equation in the plane, cf.\ \cite{CourantHilbert1962}), and its
polynomial solutions of any bounded degree are found by comparing
coefficients: they form the kernel of an explicit linear map between
finite-dimensional spaces of polynomials. Each solution $(\rho,\lambda)$ then
determines a surface: $X_1$ is recovered from the displayed gradient ---
whose compatibility is exactly \eqref{eq:wave}, which is \emph{why} the
equation governs the branch --- and $X_2,X_3$ follow by the further
integrations. For the choice $\rho=w_1$ used below no ad hoc computation is
needed: \eqref{eq:wave} reads
$\lambda_{11}+2w_1\lambda_{12}+w_1^2\lambda_{22}+2\lambda_2=0$, and the
substitution $\eta=w_2-\tfrac12w_1^2$ transforms it into the \emph{heat
equation},
\[
\lambda(w_1,w_2)\;=\;\mu\bigl(w_1,\;w_2-\tfrac12w_1^2\bigr),
\qquad
\partial_{w_1}^{2}\mu+\partial_\eta\mu\;=\;0
\]
(verified by exact expansion). Its polynomial solutions are completely
described in the literature: by Rosenbloom and Widder
\cite{RosenbloomWidder1959}, they are spanned by the \emph{heat polynomials}
\[
v_k(x,t)\;=\;\sum_{j\le k/2}\frac{k!}{j!\,(k-2j)!}\,x^{\,k-2j}\,t^{\,j},
\]
one in each degree $k$. Hence \eqref{eq:wave} with $\rho=w_1$ has exactly one
new polynomial solution in each degree, namely
$v_k\bigl(w_1,\tfrac12w_1^2-w_2\bigr)$; normalized to be monic in $w_1$, the
first six are
\[
1,\quad w_1,\quad w_1^2-w_2,\quad w_1^3-\tfrac32w_1w_2,\quad
w_1^4-\tfrac{12}5w_1^2w_2+\tfrac65w_2^2,\quad
w_1^5-\tfrac{40}{13}w_1^3w_2+\tfrac{30}{13}w_1w_2^2 .
\]
The affine $\lambda$ chosen below is a combination of the first three; every
further heat polynomial is a candidate surface for the construction, whose
fiber geometry is unexplored.

(Under the normalization $X_1=w_1$
the branch is rigid --- only a two-parameter family survives, and its twist is
obstructed; the freedom of a general $X_1$ is essential.) Taking $\rho=w_1$ and
$\lambda=\nu_1+\nu_2w_1+\nu_3(w_2-w_1^2)$, which solves \eqref{eq:wave}, and
integrating, one obtains a three-parameter family of quintic surfaces, linear
in $(\nu_1,\nu_2,\nu_3)$.
The five divisibility constants of the twist (with scalings
$(w_1,w_2)=(\gamma u_1,\gamma^2u_2)$, monomials $(v_1,v_2,v_3)=(xy,x^2z,x^3t)$
of $x$-degrees $(1,2,3)$, and twist exponents $(1,2,3)$, so that
$\det=\;C^{-6}\cdot\gamma^3\cdot\gamma^3\cdot x^6\cdot\det(\text{stage})$)
now form a \emph{solvable linear system}: at the base point
$(\gamma_0,u_1^0,u_2^0)=(1,1,1)$ one finds
$(\nu_1,\nu_2,\nu_3)=\bigl(\tfrac{51}{29},-\tfrac{60}{29},-\tfrac{240}{29}\bigr)$,
with surface
\begin{align*}
X_1&=\tfrac{160}{29}w_1^3-\tfrac{60}{29}w_1^2-\tfrac{240}{29}w_1w_2+\tfrac{51}{29}w_1+\tfrac{60}{29}w_2,\\
X_2&=\tfrac{60}{29}w_1^4-\tfrac{20}{29}w_1^3-\tfrac{120}{29}w_2^2+\tfrac{51}{29}w_2,\\
X_3&=-\tfrac{48}{29}w_1^5+\tfrac{15}{29}w_1^4+\tfrac{240}{29}w_1^3w_2-\tfrac{17}{29}w_1^3-\tfrac{60}{29}w_1^2w_2-\tfrac{240}{29}w_1w_2^2+\tfrac{51}{29}w_1w_2+\tfrac{30}{29}w_2^2 .
\end{align*}
The residual quadratic side condition is absorbed by a \emph{unipotent nonlinear
stage}: with $v=(v_1,v_2,v_3)=(xy,x^2z,x^3t)$, set
\[
\gamma=1+\tfrac{20}{29}v_1,\qquad
u_1=1-\tfrac{49}{29}v_1+8v_2,\qquad
u_2=1-\tfrac{69}{29}v_1+9v_2+v_3+\tfrac{4197}{1682}\,v_1^2 ,
\]
where the $v_1$-column lies in $\ker(\nabla E_2)\cap\ker(\nabla E_3)$, the
$v_2$-column in $\ker(\nabla E_3)$, and the coefficient $\tau=\tfrac{4197}{1682}$
kills the $v_1^2$-obstruction while leaving the stage Jacobian constant
($=\tfrac{160}{29}$), because the corresponding cofactor vanishes by design.
With $w_1=\gamma u_1$, $w_2=\gamma^2u_2$, $C=\gamma x$ and
$S_i=X_i(w_1,w_2)+\gamma\Delta_i(w_1,w_2)$, define
\[
\boxed{\;F_5\;=\;\Bigl(C,\;\;\frac{S_1}{C},\;\;\frac{S_2}{C^2},\;\;\frac{S_3}{C^3}\Bigr)^{\mathsf T}.\;}
\]
As a composition diagram,
\[
\begin{pmatrix}x\\ y\\ z\\ t\end{pmatrix}
\xrightarrow{\ \text{monomial}\ }
\begin{pmatrix}x\\ xy\\ x^2z\\ x^3t\end{pmatrix}
\xrightarrow{\ \text{stage}\ }
\begin{pmatrix}x\\ \gamma\\ u_1\\ u_2\end{pmatrix}
\xrightarrow{\ \substack{w_1=\gamma u_1\\[1pt] w_2=\gamma^2u_2}\ }
\begin{pmatrix}x\\ \gamma\\ w_1\\ w_2\end{pmatrix}
\xrightarrow{\ \text{padded sweep}\ }
\begin{pmatrix}\gamma x\\ S_1\\ S_2\\ S_3\end{pmatrix}
\xrightarrow{\ \text{twist}\ }
\begin{pmatrix}C\\ S_1/C\\ S_2/C^2\\ S_3/C^3\end{pmatrix},
\]
in which the second arrow is now the unipotent nonlinear stage displayed above
(a polynomial automorphism of the coordinates $(\gamma,u_1,u_2)$ over $x$)
rather than a plain affine map.

\begin{theorem}\label{thm:F5}
$F_5$ is a polynomial map $\C^4\to\C^4$ with components of degrees $3,12,14,16$;
its Jacobian determinant is identically $\tfrac{160}{29}$; and its generic fiber
consists of exactly $10$ points.
\end{theorem}

Consequently $F_5$ is a four-dimensional
non-injective Keller map arising from the $M$-branch, of geometric degree ten,
not equivalent to $\Phi\times\mathrm{id}$ for any three-dimensional Keller map
$\Phi$ of geometric degree $\neq10$.

\begin{proof}[Proof sketch]
The divisibilities and the constancy of the determinant are the chain
factorization above, confirmed by exact evaluation at random rational points.
For the fibers we run the same elimination scheme as for $F_4$: let
$v_1\neq0$ and set $(Y_1,Y_2,Y_3)=(v_1v_2,\,v_1^2v_3,\,v_1^3v_4)$. Eliminating
$\gamma=Y_1-X_1(w_1,w_2)$ leaves the two equations
\[
X_2+(Y_1-X_1)\,w_1\;=\;Y_2,\qquad X_3+(Y_1-X_1)\,w_2\;=\;Y_3
\]
in $(w_1,w_2)$, whose resultant with respect to $w_2$ is, after the harmless
rescaling by $-29^4/30$, the degree-ten polynomial
{\small
\begin{align*}
R(w_1;Y)={}&20480w_1^{10}-51200w_1^{9}+57600w_1^{8}+\bigl(222720Y_1-126720\bigr)w_1^{7}\\
&+\bigl(-204160Y_1-742400Y_2+234480\bigr)w_1^{6}
+\bigl(-119712Y_1+890880Y_2+890880Y_3-155448\bigr)w_1^{5}\\
&+\bigl(201840Y_1^2+59160Y_1-417600Y_2-1113600Y_3+65025\bigr)w_1^{4}\\
&+\bigl(-67280Y_1^2-538240Y_1Y_2+41412Y_1+306240Y_2+556800Y_3-44217\bigr)w_1^{3}\\
&+\bigl(-146334Y_1^2+201840Y_1Y_2+807360Y_1Y_3+75429Y_1-88740Y_2-459360Y_3\bigr)w_1^{2}\\
&+\bigl(97556Y_1^3-42891Y_1^2-403680Y_1Y_3+177480Y_3\bigr)w_1\\
&-97556Y_1^2Y_2+42891Y_1Y_2+171564Y_1Y_3-25230Y_2^2+201840Y_2Y_3-403680Y_3^2-75429Y_3 ,
\end{align*}
}
again with \emph{constant} leading coefficient, so the number of roots never
drops and, as for $F_4$, all fiber degenerations pass through $\gamma=0$. At
the witness target
$v=(\tfrac12,\tfrac23,-\tfrac13,1)$, exact univariate gcd certificates show
that $R$ is squarefree; that $R$ is coprime to the $w_2$-resultant of the
first equation with $\gamma=0$, so every root has $\gamma\neq0$; and that $R$
is coprime to the first subresultant of the pair, so each root $w_1$ carries a
unique $w_2$. Since $x=v_1/\gamma$ and the triangular stage recovers $(y,z,t)$
rationally wherever $\gamma\neq0$, that fiber has exactly $10$ points, and by
openness of these conditions the generic count is $10$. (The counts were also
confirmed by Newton iteration at multiple targets.)
\end{proof}

Written out in full,
$F_5=(F_{5,1},F_{5,2},F_{5,3},F_{5,4})^{\mathsf T}$ with
{\small
\begin{align*}
F_{5,1} ={}& \tfrac{20}{29}x^2y+x\\[2pt]
F_{5,2} ={}& \tfrac{32768000}{24389}x^7y^2z^3-\tfrac{602112000}{707281}x^6y^3z^2-\tfrac{768000}{24389}x^6y^2zt+\tfrac{3276800}{841}x^6yz^3+\tfrac{2076288000}{20511149}x^5y^4z\\
&+\tfrac{4704000}{707281}x^5y^3t-\tfrac{54835200}{24389}x^5y^2z^2+\tfrac{80752000}{20511149}x^4y^5-\tfrac{76800}{841}x^5yzt+\tfrac{81920}{29}x^5z^3+\tfrac{152428800}{707281}x^4y^3z\\
&+\tfrac{374400}{24389}x^4y^2t-\tfrac{1044480}{841}x^4yz^2+\tfrac{169140800}{20511149}x^3y^4-\tfrac{1920}{29}x^4zt+\tfrac{515520}{24389}x^3y^2z+\tfrac{3360}{841}x^3yt+\tfrac{9600}{29}x^3z^2\\
&+\tfrac{907520}{707281}x^2y^3-\tfrac{108960}{841}x^2yz-\tfrac{180}{29}x^2t+\tfrac{14050}{24389}xy^2-\tfrac{252}{29}xz+y\\[2pt]
F_{5,3} ={}& \tfrac{98304000}{24389}x^8y^2z^4-\tfrac{2408448000}{707281}x^7y^3z^3+\tfrac{9830400}{841}x^7yz^4+\tfrac{22127616000}{20511149}x^6y^4z^2-\tfrac{191692800}{24389}x^6y^2z^3\\
&-\tfrac{90354432000}{594823321}x^5y^5z-\tfrac{48000}{24389}x^6y^2t^2+\tfrac{245760}{29}x^6z^4-\tfrac{201456000}{20511149}x^5y^4t+\tfrac{45158400}{24389}x^5y^3z^2\\
&-\tfrac{73022484000}{17249876309}x^4y^6-\tfrac{864000}{24389}x^5y^2zt-\tfrac{1310720}{841}x^5yz^3-\tfrac{5316643200}{20511149}x^4y^4z-\tfrac{4800}{841}x^5yt^2-\tfrac{13521600}{707281}x^4y^3t\\
&-\tfrac{25906560}{24389}x^4y^2z^2-\tfrac{4696088400}{594823321}x^3y^5-\tfrac{86400}{841}x^4yzt+\tfrac{112640}{29}x^4z^3+\tfrac{69664320}{707281}x^3y^3z-\tfrac{120}{29}x^4t^2\\
&+\tfrac{62760}{24389}x^3y^2t-\tfrac{1613760}{841}x^3yz^2+\tfrac{40879390}{20511149}x^2y^4-\tfrac{2160}{29}x^3zt+\tfrac{3661320}{24389}x^2y^2z+\tfrac{240}{29}x^2yt+\tfrac{9480}{29}x^2z^2\\
&-\tfrac{516820}{707281}xy^3-\tfrac{105360}{841}xyz-\tfrac{189}{29}xt-\tfrac{1437}{1682}y^2-z\\[2pt]
F_{5,4} ={}& -\tfrac{629145600}{24389}x^9y^2z^5+\tfrac{19267584000}{707281}x^8y^3z^4+\tfrac{49152000}{24389}x^8y^2z^3t-\tfrac{62914560}{841}x^8yz^5-\tfrac{132882432000}{20511149}x^7y^4z^3\\
&-\tfrac{903168000}{707281}x^7y^3z^2t+\tfrac{1975910400}{24389}x^7y^2z^4-\tfrac{15504384000}{20511149}x^6y^5z^2-\tfrac{768000}{24389}x^7y^2zt^2+\tfrac{4915200}{841}x^7yz^3t\\
&-\tfrac{1572864}{29}x^7z^5+\tfrac{2308608000}{20511149}x^6y^4zt-\tfrac{15174451200}{707281}x^6y^3z^3+\tfrac{3799290048000}{17249876309}x^5y^6z+\tfrac{4704000}{707281}x^6y^3t^2\\
&-\tfrac{85708800}{24389}x^6y^2z^2t+\tfrac{54312960}{841}x^6yz^4+\tfrac{8448384000}{594823321}x^5y^5t-\tfrac{11697561600}{20511149}x^5y^4z^2+\tfrac{2437443220800}{500246412961}x^4y^7\\
&-\tfrac{76800}{841}x^6yzt^2+\tfrac{122880}{29}x^6z^3t+\tfrac{195724800}{707281}x^5y^3zt-\tfrac{496404480}{24389}x^5y^2z^3+\tfrac{184263724800}{594823321}x^4y^5z\\
&+\tfrac{374400}{24389}x^5y^2t^2-\tfrac{1873920}{841}x^5yz^2t+\tfrac{184320}{29}x^5z^4+\tfrac{484262400}{20511149}x^4y^4t+\tfrac{1251989760}{707281}x^4y^3z^2\\
&+\tfrac{135717082080}{17249876309}x^3y^6-\tfrac{1920}{29}x^5zt^2+\tfrac{2238720}{24389}x^4y^2zt-\tfrac{844800}{841}x^4yz^3-\tfrac{4465704480}{20511149}x^3y^4z+\tfrac{2760}{841}x^4yt^2\\
&+\tfrac{7680}{29}x^4z^2t-\tfrac{6755640}{707281}x^3y^3t-\tfrac{9897600}{24389}x^3y^2z^2-\tfrac{2476692498}{594823321}x^2y^5-\tfrac{145200}{841}x^3yzt+\tfrac{123776}{29}x^3z^3\\
&-\tfrac{61829880}{707281}x^2y^3z-\tfrac{210}{29}x^3t^2-\tfrac{126390}{24389}x^2y^2t-\tfrac{1938744}{841}x^2yz^2+\tfrac{128225745}{41022298}xy^4-\tfrac{2412}{29}x^2zt\\
&+\tfrac{5740350}{24389}xy^2z+\tfrac{11001}{841}xyt+\tfrac{9318}{29}xz^2+\tfrac{1629343}{1414562}y^3-\tfrac{99879}{841}yz-\tfrac{160}{29}t
\end{align*}
}

\subsection{The case $n=5$: mixed direction fields}\label{subsec:n5}

In dimension five the parameters are $w=(w_1,w_2,w_3)$, the hypersurface
$X\in\C[w]^4$ is a parametrized threefold in $\C^4$, the matrices
$P,\widetilde P$ are $3\times3$, and there are three potential branches
$L_1,L_2,L_3$. Three natural unimodular direction fields with first entry $1$
present themselves, and the general principle of \S\ref{subsec:forms} --- the
rank of $\widetilde P$ bounds the highest branch --- sorts them immediately.
Throughout we use the syzygy matrix $N$ with rows
$-\delta_{i+1}e_1^{\mathsf T}+e_{i+1}^{\mathsf T}$ and the potentials
$\Gamma=(G_2,G_3,G_4)$ of \S\ref{subsec:forms}, so that
$\widetilde P=J(\widehat\Delta)$ and the branch analysis lives in the pencil
$J(\Gamma)+\mu\widetilde P$, $\mu=X_1+\gamma$.

\begin{itemize}
\item $\Delta=(1,w_1,w_1^2,w_1^3)^{\mathsf T}$: the field depends on $w_1$ alone, so
$J(\widehat\Delta)$ has only its first column nonzero and
$\operatorname{rank}\widetilde P=1$: $L_2\equiv L_3\equiv0$ automatically, only
the bottom branch is available, and the analysis is parallel to
Specialization~I --- one sweeps a one-parameter direction field along a
threefold, using the constant-volume frame $(\Delta,\Delta',\Delta'',\Delta''')$
of its partial derivatives ($\det=12$). In \S\ref{subsec:curvefam} we solve this case
completely, in every dimension at once, and extract a second five-dimensional
counterexample from it.
\item $\Delta=(1,w_1,w_2,w_3)^{\mathsf T}$: the tautological field, with
$\widetilde P=J(w_1,w_2,w_3)^{\mathsf T}=I_3$, so $L_3\equiv1$: the top branch
is open, $\det J(F_0)=\gamma^4$, and the normalization $L_1\equiv L_2\equiv0$
generalizes the linear wave-type equation \eqref{eq:wave} of
Specialization~II; the reduction through $\rho$-multipliers goes through
verbatim, one dimension up.
\item $\Delta=(1,w_1,w_1^2,w_2)^{\mathsf T}$: the first genuinely \emph{mixed} field ---
neither curve-like nor tautological. (Any field
$(1,w_1,w_1^2,\ell(w_2,w_3))^{\mathsf T}$ with $\ell$ a nonconstant linear
form, e.g.\ $w_2+w_3$, reduces to this one by an affine change of the
parameters $(w_2,w_3)$, which only rescales $\det J(F_0)$.) This is the case
we now work out; it looks very different, and indeed behaves differently: its
middle branch turns out to be rigid on a natural locus.
\end{itemize}

\medskip\noindent\textbf{The pencil for $\Delta=(1,w_1,w_1^2,w_2)^{\mathsf T}$.}
Here
\[
\partial_3\Delta=0,\qquad
\widehat\Delta=(w_1,w_1^2,w_2)^{\mathsf T},\qquad
\widetilde P=J(\widehat\Delta)=
\begin{pmatrix}1&0&0\\ 2w_1&0&0\\ 0&1&0\end{pmatrix},
\]
of rank $2$: the top branch dies, $L_3\equiv0$, and both the bottom and the
middle branch are a priori available.

The frame viewpoint of \S\ref{subsec:n4} applies here as well, provided one
takes \emph{all} independent partial derivatives of $\Delta$, including the
second-order one: $(\Delta,\partial_1\Delta,\partial_1^2\Delta,\partial_2\Delta)$
is a frame of constant volume $2$, with
\[
N\Delta=0,\qquad
N\partial_1\Delta=(1,\,2w_1,\,0)^{\mathsf T},\qquad
N\partial_1^2\Delta=(0,\,2,\,0)^{\mathsf T},\qquad
N\partial_2\Delta=(0,\,0,\,1)^{\mathsf T},
\]
and expanding
$\partial_jX=c_{0j}\Delta+c_{1j}\partial_1\Delta+c_{2j}\partial_1^2\Delta
+c_{3j}\partial_2\Delta$ forces $c_{0j}=\partial_jX_1$, as always, and gives
the columns of $P=N\,J(X)$ as $(c_{1j},\,2w_1c_{1j}+2c_{2j},\,c_{3j})^{\mathsf T}$
(all verified by expansion). The difference with \S\ref{subsec:n4} lies in
what must be done with these first-order data: the branch equations below have
to be \emph{integrated} --- twice on the middle branch --- so it is more
efficient to work from the start with the integrated objects, the potentials
of \S\ref{subsec:forms}, of whose derivatives the frame coefficients are
triangular combinations: $P=J(\Gamma)+X_1\widetilde P$. With the potentials
\[
G_2=X_2-w_1X_1,\qquad G_3=X_3-w_1^2X_1,\qquad G_4=X_4-w_2X_1 ,
\]
the pencil determinant is quadratic in $\mu$,
\[
\det\bigl(J(\Gamma)+\mu\widetilde P\bigr)\;=\;D_0+\mu D_1+\mu^2D_2 ,
\]
with coefficients (writing $H:=G_3-2w_1G_2$)
\begin{align*}
D_2&=-\,\partial_3H,\\
D_1&=(\partial_2H)\,\partial_3G_4-(\partial_3H)\,\partial_2G_4
-\bigl(\partial_1G_2\,\partial_3G_3-\partial_1G_3\,\partial_3G_2\bigr),\\
D_0&=\det J(G_2,G_3,G_4),
\end{align*}
and the padded sweep of Lemma~\ref{lem:JL} satisfies the exact identity
\begin{equation}\label{eq:pencil5}
\det J(F_0)\;=\;\gamma\,\Bigl[D_0+(X_1+\gamma)D_1+(X_1+\gamma)^2D_2\Bigr],
\end{equation}
verified by expansion for generic data. Comparing coefficients of $\gamma$,
the two branches become:
\begin{itemize}
\item \emph{bottom ($L_1$-)branch}: $D_2\equiv0$, $D_1\equiv c\neq0$, and
$X_1=-D_0/c$; then $\det J(F_0)=c\,\gamma^2$;
\item \emph{middle ($L_2$-)branch}: $D_2\equiv c\neq0$, $X_1=-D_1/(2c)$, and
the \emph{master equation}
\begin{equation}\label{eq:master5}
4c\,D_0\;=\;D_1^{\,2} ;
\end{equation}
then $\det J(F_0)=c\,\gamma^3$.
\end{itemize}
The two branches behave very differently --- the bottom one is solved in
closed form and yields our counterexample, while the middle one turns out to
be rigid on a natural locus --- and we treat them in turn.

\subsubsection{The bottom branch: closed-form solution and the map $F_6$}
\label{subsubsec:bottom5}

Here $D_2=-\partial_3H\equiv0$ says that $H=G_3-2w_1G_2$ does not involve
$w_3$:
\[
G_3\;=\;2w_1G_2+\varphi(w_1,w_2),
\qquad\text{$\varphi$ a free bivariate polynomial.}
\]
The remaining condition $D_1\equiv c$ is then solved in closed form.
Substituting this expression for $G_3$ into the displayed formula for $D_1$
--- so that $\partial_3H=0$, $\partial_2H=\partial_2\varphi$,
$\partial_3G_3=2w_1\,\partial_3G_2$ and
$\partial_1G_3=2G_2+2w_1\partial_1G_2+\partial_1\varphi$ --- two of the four
terms cancel and $D_1$ collapses to
\[
D_1\;=\;(\partial_2\varphi)\,\partial_3G_4
\;+\;\bigl(2G_2+\partial_1\varphi\bigr)\,\partial_3G_2 :
\]
a single first-order equation for $G_4$, with $G_2$ and $\varphi$ free. Taking
$\varphi=\tau w_2$ ($\tau\neq0$ constant --- the choice that makes the
left-hand side a total $\partial_3$-derivative, and the one compatible with
the weights of the twist below), the condition $D_1\equiv c$ reads
$\partial_3\bigl(\tau G_4+G_2^{\,2}-c\,w_3\bigr)=0$ and integrates at once.
The complete solution of the branch is therefore
\[
G_2\in\C[w_1,w_2,w_3]\ \text{and}\ \psi\in\C[w_1,w_2]\ \text{arbitrary},\qquad
G_3=2w_1G_2+\tau w_2,\qquad
G_4=\frac{c\,w_3-G_2^{\,2}}{\tau}+\psi ,
\]
followed by $X_1=-D_0/c$ and $X_{i+1}=G_{i+1}+\delta_{i+1}X_1$, which give
$\det J(F_0)=c\gamma^2$ for \emph{every} such choice: the branch is
\emph{flexible} --- an arbitrary trivariate $G_2$ generates a solution.



We now choose data with genuine, nonlinear dependence on all three
parameters:
\[
\tau=c=1,\qquad
G_2=w_2+w_1w_3+w_1^2-w_1^3,\qquad
\varphi=w_2,\qquad
\psi=2w_1^4-2w_1^5 ,
\]
so that $G_3=2w_1G_2+w_2$, $G_4=w_3-G_2^{\,2}+2w_1^4-2w_1^5$, and
$X_1=-D_0$. Explicitly,
\begin{align*}
X_1&=-w_3+2w_2-2w_1+2w_1w_3+5w_1^2-2w_1^3+8w_1^4-10w_1^5,\\
X_2&=w_2+2w_1w_2-w_1^2+2w_1^2w_3+4w_1^3-2w_1^4+8w_1^5-10w_1^6,\\
X_3&=w_2+2w_1w_2+w_1^2w_3+2w_1^2w_2+2w_1^3w_3+3w_1^4-2w_1^5+8w_1^6-10w_1^7,\\
X_4&=w_3-w_2w_3+w_2^2-2w_1w_2+3w_1^2w_2-w_1^2w_3^2-2w_1^3w_3+w_1^4+2w_1^4w_3+8w_1^4w_2-10w_1^5w_2-w_1^6 ,
\end{align*}
of degrees $5,6,7,6$. The sweep $S=X+\gamma\Delta$ satisfies the polynomial
identity $\det J(S)=\gamma$ (variables ordered $(\gamma,w_1,w_2,w_3)$), the
threefold $X$ has $\operatorname{rank}J(X)=3$ generically, and $w_3$ enters $X$
quadratically: in particular $X$ does not factor through two parameters, which
rules out the suspensions that dominate the middle branch
(Theorem~\ref{thm:collapse5} below). The pure-$w_1$ tails
$t_2=w_1^2-w_1^3\subset G_2$ and $t_4=\psi=2w_1^4-2w_1^5$ serve a second
purpose: they arrange the base point that the twist below needs. Indeed, since
$S_1=X_1+\gamma$ and $S_{i+1}=G_{i+1}+\delta_{i+1}(X_1+\gamma)$, the base
condition $S(1,1,0,0)=0$ is equivalent to the four scalar conditions
\[
G_2(1,0,0)=G_3(1,0,0)=G_4(1,0,0)=0,\qquad
X_1(1,0,0)=-1\ \ \text{(i.e.\ }D_0(1,0,0)=c=1\text{)}.
\]
Because $G_3(1,0,0)=2\,G_2(1,0,0)$ and $G_4(1,0,0)=-G_2(1,0,0)^2+\psi(1,0)$,
the first three reduce to $t_2(1)=0$ and $t_4(1)=0$, which the displayed tails
satisfy; the last, $D_0(1,0,0)=1$, is a linear condition on the leading tail
coefficients and holds by direct evaluation (all verified exactly).

For the twist, assign the weights $(1,3,4)$ to $(w_1,w_2,w_3)$. Every monomial of
$X_i$ has weight $\ge i$ against the twist exponents $(1,2,3,4)$, so the
scalings
\[
w_1=\gamma u_1,\qquad w_2=\gamma^3u_2,\qquad w_3=\gamma^4u_3
\]
make $E_i:=S_i/\gamma^{\,i}$ a polynomial in $(\gamma,u_1,u_2,u_3)$ for
$i=1,2,3,4$, with $E_i(1,1,0,0)=0$ at the base point. The gradients there are
\[
\nabla E_2=(-16,-17,3,2),\qquad
\nabla E_3=(-17,-18,5,3),\qquad
\nabla E_4=(-2,-2,0,1).
\]
On $\C^5\ni(x,y,z_1,z_2,z_3)$ --- we name the last three source coordinates
$z_1,z_2,z_3$ to keep them clear of the sweep parameters --- put $v=(v_1,v_2,v_3,v_4)=(xy,\,x^2z_1,\,x^3z_2,\,x^4z_3)$
(of $x$-degrees $1,2,3,4$) and define the \emph{stage}, solving the
hierarchy \eqref{eq:side},
\[
\begin{aligned}
\gamma&=1-29v_1+999v_1^2+355v_1v_2-41553v_1^3+v_4,\\
u_1&=1+27v_1-5v_2,\\
u_2&=v_1-12v_2+\tfrac{2128}{5}v_1^2+v_3,\\
u_3&=-4v_1-10v_2 .
\end{aligned}
\]
Here the $v_1$-column $(-29,27,1,-4)$ spans the common kernel
$\ker\nabla E_2\cap\ker\nabla E_3\cap\ker\nabla E_4$, the $v_2$-column
$(0,-5,-12,-10)$ lies in $\ker\nabla E_3\cap\ker\nabla E_4$, the $v_3$- and
$v_4$-columns are $e_{u_2}$ and $e_\gamma$ with
$\nabla E_4\cdot e_{u_2}=0$, and the three correction vectors (the
coefficients of $v_1^2$, $v_1v_2$, $v_1^3$) are chosen inside
$\mathrm{span}(e_{u_2},e_\gamma)$ --- the span of the last two columns --- to
kill the obstructing Taylor coefficients of $E_3$ and $E_4$ along the
$v_1$-direction. Two consequences: first, $E_i$ composed with the stage is
divisible by $x^{\,i}$; second, because the corrections lie in the span of two
columns of the linear part, the stage is an \emph{elementary polynomial
automorphism} of $\C^4$ with constant Jacobian determinant $-290$ (each
correction's contribution to the Jacobian is a determinant with a repeated
column). With $w_1=\gamma u_1$, $w_2=\gamma^3u_2$, $w_3=\gamma^4u_3$ substituted
throughout, the sweep components are
\[
S_i\;=\;X_i(w_1,w_2,w_3)+\gamma\,\Delta_i(w_1,w_2)\quad(i=1,2,3,4),
\quad\text{i.e.}\quad
\begin{aligned}
S_1&=X_1+\gamma, & S_2&=X_2+\gamma\,w_1,\\
S_3&=X_3+\gamma\,w_1^2, & S_4&=X_4+\gamma\,w_2 ,
\end{aligned}
\]
with $X_1,\dots,X_4$ the polynomials displayed above. With $C=\gamma x$, define
\[
\boxed{\;F_6\;=\;\Bigl(C,\;\;\frac{S_1}{C},\;\;\frac{S_2}{C^2},\;\;
\frac{S_3}{C^3},\;\;\frac{S_4}{C^4}\Bigr)^{\mathsf T}\;:\;\C^5\to\C^5 .\;}
\]
As a composition diagram, in the format of \S\ref{subsec:twist},
\begin{multline*}
\begin{pmatrix}x\\ y\\ z_1\\ z_2\\ z_3\end{pmatrix}
\;\xrightarrow{\ \text{monomial}\ }\;
\begin{pmatrix}x\\ xy\\ x^2z_1\\ x^3z_2\\ x^4z_3\end{pmatrix}
\;\xrightarrow{\ \text{stage}\ }\;
\begin{pmatrix}x\\ \gamma\\ u_1\\ u_2\\ u_3\end{pmatrix}
\;\xrightarrow{\ \substack{w_1=\gamma u_1\\[1pt] w_2=\gamma^3u_2\\[1pt] w_3=\gamma^4u_3}\ }\;
\begin{pmatrix}x\\ \gamma\\ w_1\\ w_2\\ w_3\end{pmatrix}\\
\xrightarrow{\ \text{padded sweep}\ }\;
\begin{pmatrix}\gamma x\\ S_1\\ S_2\\ S_3\\ S_4\end{pmatrix}
\;\xrightarrow{\ \text{twist}\ }\;
\begin{pmatrix}C\\ S_1/C\\ S_2/C^2\\ S_3/C^3\\ S_4/C^4\end{pmatrix},
\end{multline*}
where the stage is the elementary polynomial automorphism displayed above
(over $x$), of constant Jacobian determinant $-290$; the Jacobian determinants
of the five arrows multiply to the constant in the chain factorization of the
proof of Theorem~\ref{thm:F6}.

\begin{theorem}\label{thm:F6}
$F_6$ is a polynomial map $\C^5\to\C^5$ with components of degrees
$7,38,40,42,44$; its Jacobian determinant is identically $-290$; and its
generic fiber consists of exactly $6$ points.
\end{theorem}

The proof below establishes two further properties: $F_6$ fixes the axis
$\{y=z_1=z_2=z_3=0\}$ pointwise, and for every $C_0\neq0$ the fiber over
$(C_0,0,0,0,0)$ consists of exactly $4$ points --- the visible one and three
escape companions. Consequently $F_6$ is a non-injective Keller map in
dimension five, of geometric degree six, not equivalent to
$\Phi\times\mathrm{id}$ for any four-dimensional Keller map $\Phi$ of geometric
degree $\neq6$, since geometric degree is invariant under composition with
polynomial automorphisms and multiplicative on products.

\begin{proof}[Proof sketch]
The Jacobian statement is the chain factorization
\[
\det J(F_6)\;=\;
\underbrace{x^{10}}_{\text{monomials}}\cdot
\underbrace{(-290)}_{\text{stage}}\cdot
\underbrace{\gamma^{8}}_{\text{scalings }\gamma\cdot\gamma^3\cdot\gamma^4}\cdot
\underbrace{\gamma^{2}}_{\det J(F_0)=1\cdot\gamma^2}\cdot
\underbrace{(\gamma x)^{-10}}_{\text{twist exponents }1+2+3+4}
\;=\;-290 ,
\]
together with the divisibility statements above, all verified exactly (and the
constancy confirmed at $17$ random rational points). For the fibers, let
$Y\in\C^4$ be a sweep target. The system $S(\gamma,w)=Y$ \emph{triangularizes}:
$S_1=Y_1$ gives $\gamma=Y_1-X_1$; substituting into
$S_{i+1}=G_{i+1}+\delta_{i+1}(X_1+\gamma)$ turns the remaining three equations
into the potential equations
\[
G_2(w)=Y_2-w_1Y_1,\qquad G_3(w)=Y_3-w_1^2Y_1,\qquad G_4(w)=Y_4-w_2Y_1 ,
\]
By our choice of the data, $G_3=2w_1G_2+w_2$ and
$G_4=w_3-G_2^{\,2}+\psi(w_1)$ with $\psi=2w_1^4-2w_1^5$; hence
\[
w_2\;=\;G_3-2w_1G_2\;=\;\bigl(Y_3-w_1^2Y_1\bigr)-2w_1\bigl(Y_2-w_1Y_1\bigr)
\;=\;Y_3-2w_1Y_2+w_1^2Y_1
\]
and, inserting $G_2=Y_2-w_1Y_1$ and this value of $w_2$ (the mixed terms
cancel identically),
\[
w_3\;=\;Y_4-w_2Y_1+G_2^{\,2}-\psi(w_1)
\;=\;Y_4+Y_2^{\,2}-Y_1Y_3+2w_1^5-2w_1^4 .
\]
Substituting both expressions into the one equation not yet used,
$G_2\bigl(w_1,w_2(w_1),w_3(w_1)\bigr)=Y_2-w_1Y_1$, leaves the univariate \emph{fiber polynomial}
\[
R_Y(w_1)\;=\;2w_1^6-2w_1^5-w_1^3+(Y_1+1)\,w_1^2
+\bigl(Y_4+Y_2^{\,2}-Y_1Y_3-2Y_2+Y_1\bigr)\,w_1+\bigl(Y_3-Y_2\bigr)\;=\;0 ,
\]
of degree exactly $6$ with \emph{constant} leading coefficient $2$. Since the
stage is a polynomial automorphism and the monomial conjugation is bijective
wherever $x\neq0$, the fiber of $F_6$ over a target with first coordinate
$C_0\neq0$ is in bijection with the roots of $R_Y$ having $\gamma=Y_1-X_1\neq0$,
via $x=C_0/\gamma$. For random rational targets $R_Y$ is squarefree with all
six roots satisfying $\gamma\neq0$ (verified exactly for the squarefreeness and
numerically for the roots), giving generic fiber size $6$. Over
$Y=0$ the fiber polynomial factors,
\[
R_0(w_1)\;=\;w_1^2\,(w_1-1)\,(2w_1^3-1),
\]
and $\gamma=-X_1$ vanishes at the double root $w_1=0$ (which is therefore
deleted by the twist) while $\gamma=1$ at $w_1=1$ and
$\gamma=3w_1^2-3w_1=3w_1(w_1-1)$ at the three roots of $2w_1^3=1$; since
$3w_1(w_1-1)$ vanishes only at $w_1\in\{0,1\}$, which are not roots of
$2w_1^3-1$, all four remaining points survive, and the fiber over
$(C_0,0,0,0,0)$ has exactly $4$ points --- the visible one
$(C_0,0,0,0,0)$ (note $S(1,1,0,0)=0$ makes $F_6$ fix the axis) and three
escape companions.
\end{proof}

The first component of $F_6$ is
\[
F_{6,1}\;=\;x-29x^2y+999x^3y^2-41553x^4y^3+355x^4yz_1+x^5z_3 ,
\]
and the remaining components have $342$, $421$, $507$ and $904$ terms
respectively; they are reproduced, together with all verification scripts, in
the ancillary files. Unlike $F_4$ and $F_5$, whose geometric degrees ($5$ and
$10$) were established by elimination of comparable difficulty, here the
triangular collapse of the fiber system to the sextic $R_Y$ makes the
five-dimensional example the \emph{easiest} of the three to analyze --- a
direct payoff of the potentials $(G_2,G_3,G_4)$ adapted to the mixed direction
field.

\subsubsection{The middle branch: rigidity}\label{subsubsec:middle5}

Here $D_2\equiv c$ integrates to
\[
G_3\;=\;2w_1G_2-c\,w_3+\varphi(w_1,w_2),
\qquad\text{$\varphi$ a free bivariate polynomial,}
\]
$X_1=-D_1/(2c)$ is determined, and the data $(G_2,G_4,\varphi)$ remain
constrained by the nonlinear master equation \eqref{eq:master5}: unlike its
neighbour, this branch does not integrate in closed form. Introduce the
\emph{characteristic field} and its invariant coordinate
\[
\mathcal W\;=\;c\,\partial_2+(\partial_2\varphi)\,\partial_3 ,
\qquad
\theta\;=\;w_3-\varphi(w_1,w_2)/c
\qquad(\mathcal W\theta=0,\ \ \mathcal W w_1=0).
\]
The quadratic part of \eqref{eq:master5} in the derivatives of $G_4$ is the
rank-one form $(\mathcal W G_4)^2$, so \eqref{eq:master5} is a
Hamilton--Jacobi-type equation degenerate along $\mathcal W$.

\begin{theorem}[collapse on the characteristic-invariance locus]\label{thm:collapse5}
Let $c\neq0$, let $\varphi(w_1,w_2)$ be arbitrary, and suppose the datum $G_2$ is
constant along the characteristic field, $\mathcal WG_2=0$, i.e.\
$G_2=G_2(w_1,\theta)$. Set
\[
K(w_1,\theta)\;=\;c\,\frac{\partial G_2}{\partial w_1}\Big|_{\theta}
\;+\;\frac{\partial\,(G_2^{\,2})}{\partial\theta}.
\]
Then the master equation \eqref{eq:master5} factors as a perfect square,
\[
4c\,D_0-D_1^{\,2}\;=\;-\bigl(\mathcal WG_4-K\bigr)^2 ,
\]
so it linearizes to the transport equation $\mathcal WG_4=K$, whose general
polynomial solution is $G_4=(w_2/c)\,K(w_1,\theta)+h(w_1,\theta)$ with $h$
arbitrary. For every such solution the components $X_1,X_2,X_3$ of the swept
threefold depend on $(w_1,\theta)$ alone, $X_4=h(w_1,\theta)$, and in the
parameters $(\gamma,w_1,\theta,w_2)$ the sweep is the triangular suspension
\[
S\;=\;\bigl(T(\gamma,w_1,\theta),\ \gamma w_2+h(w_1,\theta)\bigr),
\qquad
T\;=\;\widehat X(w_1,\theta)+\gamma\,(1,w_1,w_1^2)^{\mathsf T},
\qquad
\det J(T)=c\,\gamma:
\]
a four-dimensional bottom-branch sweep with direction field $(1,w_1,w_1^2)^{\mathsf T}$,
extended by one shear-type coordinate. In particular no datum with
$\mathcal WG_2=0$ produces a genuinely five-dimensional example.
\end{theorem}

\begin{proof}
All statements are direct computations from the displayed formulas for
$D_0,D_1,D_2$. With $H=-cw_3+\varphi$ one finds
$\partial_2H=\partial_2\varphi$, $\partial_3H=-c$, and, using
$\partial_2G_2=-((\partial_2\varphi)/c)\partial_\theta G_2$,
$\partial_3G_2=\partial_\theta G_2$:
the coefficient of $\partial_1G_4$ in $D_0$ is
$-(c\,\partial_2G_2+(\partial_2\varphi)\partial_3G_2)=-\mathcal WG_2=0$; the remaining
two coefficients of $D_0$ collapse to $(K/c)\,\mathcal W$, giving
$D_0=(K/c)\,\mathcal WG_4$; and $D_1=\mathcal WG_4+K$. Hence
$4cD_0-D_1^2=4K\,\mathcal WG_4-(\mathcal WG_4+K)^2=-(\mathcal WG_4-K)^2$.
Since $\mathcal W$ annihilates functions of $(w_1,\theta)$ and
$\mathcal W\bigl((w_2/c)K\bigr)=K$, the transport equation has the stated general
solution. Finally $X_1=-D_1/(2c)=-K/c$ depends on $(w_1,\theta)$ only, hence so
do $X_2=G_2+w_1X_1$ and
$X_3=G_3+w_1^2X_1=2w_1G_2+\varphi-cw_3+w_1^2X_1=2w_1G_2-c\theta+w_1^2X_1$, while
$X_4=G_4+w_2X_1=h$; and $\partial S_4/\partial w_2=\gamma$ with
$S_1,S_2,S_3$ independent of $w_2$ gives $\det J(S)=\gamma\det J(T)$, so
$\det J(T)=c\gamma$. All identities were also verified by exact expansion on
randomized data.
\end{proof}

\begin{remark}
The collapse is not confined to the locus $\mathcal WG_2=0$. The
complementary explicit families we have solved --- e.g.\ $\varphi=0$,
$G_2$ independent of $w_3$ and $G_4$ linear in $w_3$ with constant coefficient ---
also collapse, in a different mode ($w_3$ enters all components linearly through
constant coefficients, and a linear change of target coordinates splits off a
graph coordinate).
\end{remark}

\begin{problem}[rigidity of the middle branch]\label{prob:rigid5}
Does the master equation \eqref{eq:master5} admit a polynomial solution that is
not equivalent to a suspension of a four-dimensional sweep? We conjecture that
it does not. If so, this is a genuine rigidity phenomenon: for the mixed field
$\Delta=(1,w_1,w_1^2,w_2)^{\mathsf T}$ the middle escape $\det J(F_0)=c\gamma^3$ is
unrealizable by irreducibly five-dimensional data, even though it is
unobstructed at the level of branch counting.
\end{problem}

\subsection{The curve-type family: a uniform solution in every dimension}
\label{subsec:curvefam}

We now return to the first direction field of the classification above and
treat it uniformly for all $n\ge3$:
\[
\Delta\;=\;\bigl(1,\,w_1,\,w_1^2,\,\dots,\,w_1^{\,n-2}\bigr)^{\mathsf T},
\]
the rational normal curve, depending on the single parameter $w_1$. Here the
normalization system admits a \emph{complete} solution, and the solutions for
$n=3,4,5$ fall into a single pattern.

\begin{theorem}[reduction for the curve-type field]\label{thm:curvered}
Let $n\ge3$ and $\Delta=(1,w_1,\dots,w_1^{\,n-2})^{\mathsf T}$. With the
potentials $\Gamma=(G_2,\dots,G_{n-1})$ of \S\ref{subsec:forms}, set
\[
H_{i+1}\;:=\;G_{i+1}-i\,w_1^{\,i-1}G_2\qquad(2\le i\le n-2).
\]
Then $\widetilde P=\widehat\Delta{}'\,e_1^{\mathsf T}$ is of rank one, with
$\widehat\Delta{}'=(1,2w_1,\dots,(n-2)w_1^{\,n-3})^{\mathsf T}$, the pencil
is linear in $\mu$,
\[
\det\bigl(J(\Gamma)+\mu\widetilde P\bigr)\;=\;D_0+\mu D_1,
\qquad
D_0=\det J(\Gamma),
\qquad
D_1\;=\;\det J_{(w_2,\dots,w_{n-2})}\bigl(H_3,\dots,H_{n-1}\bigr),
\]
and the (unique) branch normalization $D_1\equiv c\neq0$ says precisely:
\begin{itemize}
\item[(i)] the family
$H(w_1;\,\cdot\,)=(H_3,\dots,H_{n-1})^{\mathsf T}\colon\C^{n-3}\to\C^{n-3}$ has constant
Jacobian determinant $c$ in the transverse parameters $(w_2,\dots,w_{n-2})$,
for every value of the spectator parameter $w_1$;
\item[(ii)] $G_2\in\C[w]$ is completely free.
\end{itemize}
Given such data, $X_1=-D_0/c$ and $X_{i+1}=G_{i+1}+w_1^{\,i}X_1$ yield
$\det J(F_0)=c\,\gamma^2$.
\end{theorem}

\begin{proof}
Since $\Delta$ depends on $w_1$ alone, $J(\Delta)$ has only its first column
nonzero and $\widetilde P=N\,J(\Delta)=\widehat\Delta{}'\,e_1^{\mathsf T}$;
a rank-one summand makes the pencil determinant linear
in $\mu$, with $\mu$-coefficient $D_1=\det(\widehat\Delta{}',\partial_2\Gamma,\dots,
\partial_{n-2}\Gamma)$. Subtracting $i\,w_1^{\,i-1}$ times the first row from
the $i$-th row ($2\le i\le n-2$) reduces the column $\widehat\Delta{}'$ to $e_1$; these row
operations commute with $\partial_j$ for $j\ge2$ because the multipliers
depend on $w_1$ alone, and they replace $\partial_jG_{i+1}$ by
$\partial_jH_{i+1}$. Expanding along the first column leaves the stated
$(n-3)\times(n-3)$ Jacobian. Finally $X_1=-D_0/c$ gives
$\det J(S)=D_0+(X_1+\gamma)c=c\gamma$, and padding contributes the extra
$\gamma$. All identities were verified by exact expansion on randomized data.
\end{proof}

\begin{remark}
The frame viewpoint of \S\ref{subsec:n4} is available here too: the collection
$(\Delta,\Delta',\dots,\Delta^{(n-2)})$ of \emph{all} independent partial
derivatives of the rational normal field is a frame of constant volume
$\prod_{k=0}^{n-2}k!$ ($=2$ for $n=4$ and $12$ for $n=5$, the values already
met), and expanding the columns of $J(X)$ in it --- the coefficient of
$\Delta$ being $\partial_jX_1$, as always --- reproduces the pencil above; the
row reduction in the proof is precisely the triangular change between the
frame coefficients and the derivatives of the potentials. We phrase the
theorem through the potentials because the construction ultimately
\emph{integrates} the data: $G_2$ enters the solution undifferentiated, and
the fiber theory of Proposition~\ref{prop:curvefib} below is written directly
in terms of $\Gamma$.
\end{remark}

The pattern for small $n$ is now transparent. For $n=3$ there are no
transverse parameters: the condition (i) is empty ($D_1=1$, $c=1$), $G_2(w_1)$
is the only datum, and $X=(-G_2',\,G_2-w_1G_2')^{\mathsf T}$ sweeps the envelope of the
line family $\{Y_2=w_1Y_1+G_2(w_1)\}$ --- the classical Clairaut picture, with
$G_2$ the support function; the whole of \S\ref{sec:sweep} is the case $n=3$.
For $n=4$ the condition reads $\partial_2H_3=c$, i.e.\ $H_3=cw_2+\varphi(w_1)$;
taking $G_2$ \emph{linear} in $w_2$ recovers exactly the ruled surfaces of
Specialization~I, while a nonlinear $G_2$ gives non-ruled solutions that the
ruled ansatz of \S\ref{subsec:specI} does not see. For $n=5$ the condition
asks for a $w_1$-family of \emph{unit-Jacobian pairs}
$(H_3,H_4)$ in $(w_2,w_3)$ --- the first case where (i) is a genuine
Keller-type condition on the data.

\begin{proposition}[univariate fiber theory]\label{prop:curvefib}
Assume in addition that each $H(w_1;\,\cdot\,)$ is a polynomial automorphism
of $\C^{n-3}$ whose inverse depends polynomially on $w_1$ (e.g.\ any
triangular family). Then for every target $Y\in\C^{n-1}$ the fiber system
$S(\gamma,w)=Y$ triangularizes: $\gamma=Y_1-X_1$, the potential equations
become $G_{i+1}(w)=Y_{i+1}-w_1^{\,i}Y_1$, hence
\[
H\bigl(w_1;\,w_2,\dots,w_{n-2}\bigr)\;=\;h(w_1),
\qquad
h_{i+1}\;=\;Y_{i+1}-i\,w_1^{\,i-1}Y_2+(i-1)\,w_1^{\,i}Y_1 ,
\]
so $(w_2,\dots,w_{n-2})=H^{-1}(w_1;h(w_1))$ are polynomial functions of $w_1$,
and the last remaining equation is the univariate \emph{fiber polynomial}
\[
R_Y(w_1)\;:=\;G_2\bigl(w_1,\,H^{-1}(w_1;h(w_1))\bigr)-\bigl(Y_2-w_1Y_1\bigr)
\;=\;0 .
\]
The fibers of the twisted map over targets with nonzero first coordinate are
in bijection with the roots of $R_Y$ at which $\gamma\neq0$. For $n=3$, $R_Y$
is the tangency polynomial of \S\ref{sec:sweep}.
\end{proposition}

\begin{remark}[counterexamples propagate along $n\mapsto n+3$]
Condition (i) of Theorem~\ref{thm:curvered} asks for a $w_1$-family of
constant-Jacobian --- Keller-type --- maps of $\C^{n-3}$. Any family of
(tame) polynomial automorphisms qualifies, and then
Proposition~\ref{prop:curvefib} applies. But for $n\ge6$ one may instead take
$H$ to be a \emph{constant} family equal to a Keller counterexample in $n-3$
variables: Alp\"oge's $F$ for $n=6$, the maps $F_4$ or $F_5$ for $n=7$, $F_6$
or $F_7$ for $n=8$, and so on --- the construction feeds on its own output,
three dimensions at a time. (The sweep-level statement is unconditional;
carrying out the monomial twist requires the same finitely many side
conditions as in the examples above.) For such non-invertible $H$ the fiber
system is no longer triangular, and its geometry is unexplored.
\end{remark}

\medskip\noindent\textbf{An explicit five-dimensional example.} Take $n=5$,
$c=1$, and the data
\[
H_3=w_2+w_1^3+w_1^2w_3,\qquad
H_4=w_3+w_1^4,\qquad
G_2=w_2^{\,2}+w_2+w_1w_3+w_1^3 ,
\]
so that $\det J_{(w_2,w_3)}(H_3,H_4)\equiv1$, the pair
$H(w_1;\,\cdot\,)$ is a triangular automorphism for every $w_1$, and $G_2$ is
genuinely nonlinear in the transverse parameters. Then $G_3=2w_1G_2+H_3$,
$G_4=3w_1^2G_2+H_4$, $X_1=-\det J(G_2,G_3,G_4)$, and the components
$X_1,\dots,X_4$ (degrees $7,8,9,10$) are reconstructed as in
Theorem~\ref{thm:curvered}; the sweep $S=X+\gamma\Delta$ satisfies
$\det J(S)=\gamma$ and passes through the origin at the base point,
$S(1,1,0,-1)=0$, with transversality constant $D_0(1,0,-1)=1$ (the pure-$w_1$
terms of $H_3$, $H_4$, $G_2$ are tuned for exactly this). Assign the weights
$(1,3,4)$ to $(w_1,w_2,w_3)$; every monomial of $X_i$ has weight $\ge i$, so
the scalings $w_1=\gamma u_1$, $w_2=\gamma^3u_2$, $w_3=\gamma^4u_3$ make
$E_i=S_i/\gamma^{\,i}$ polynomial, with base point
$(\gamma,u_1,u_2,u_3)=(1,1,0,-1)$ and gradients
\[
\nabla E_2=(-15,0,-3,4),\qquad
\nabla E_3=(-20,3,-1,6),\qquad
\nabla E_4=(-19,8,-1,7).
\]
On $\C^5\ni(x,y,z_1,z_2,z_3)$ with $v=(xy,\,x^2z_1,\,x^3z_2,\,x^4z_3)$ the
stage is
\[
\begin{aligned}
\gamma&=1-61v_1+9012v_1^2+\tfrac{238}{19}v_1v_2-\tfrac{35823670}{19}v_1^3+8v_3+v_4,\\
u_1&=1+59v_1+v_2-5178v_1^2+19v_3,\\
u_2&=-7v_1-27v_2,\\
u_3&=-1-234v_1-5v_2 ,
\end{aligned}
\]
built by the same rules as for $F_6$: the $v_1$-column $(-61,59,-7,-234)$
spans the common kernel of the three gradients, the $v_2$-column
$(0,1,-27,-5)$ lies in $\ker\nabla E_3\cap\ker\nabla E_4$, the $v_3$- and
$v_4$-columns $(8,19,0,0)$ and $e_\gamma$ satisfy $\nabla E_4\cdot(8,19,0,0)=0$,
and all correction vectors lie in the span of the last two columns, so the
stage is an elementary polynomial automorphism of constant Jacobian
determinant $119377$. With $C=\gamma x$ and
$S_i=X_i(w_1,w_2,w_3)+\gamma\,w_1^{\,i-1}$ substituted throughout, define
\[
\boxed{\;F_7\;=\;\Bigl(C,\;\;\frac{S_1}{C},\;\;\frac{S_2}{C^2},\;\;
\frac{S_3}{C^3},\;\;\frac{S_4}{C^4}\Bigr)^{\mathsf T}\;:\;\C^5\to\C^5 .\;}
\]

\begin{theorem}\label{thm:F7}
$F_7$ is a polynomial map $\C^5\to\C^5$ with components of degrees
$7,86,89,92,95$; its Jacobian determinant is identically $119377$; and its
generic fiber consists of exactly $12$ points.
\end{theorem}

The determinant is the chain factorization
$x^{10}\cdot119377\cdot\gamma^{8}\cdot\gamma^{2}\cdot(\gamma x)^{-10}$, exactly
as for $F_6$. As before, the proof yields more: $F_7$ fixes the axis
$\{y=z_1=z_2=z_3=0\}$ pointwise, and for $C_0\neq0$ the fiber over
$(C_0,0,0,0,0)$ consists of exactly $7$ points. In particular $F_7$ is a
five-dimensional non-injective Keller map of geometric degree twelve, not
equivalent to $\Phi\times\mathrm{id}$ for any four-dimensional Keller map
$\Phi$ of geometric degree $\neq12$.

\begin{proof}[Proof sketch]
The Jacobian and divisibility statements are verified exactly as for $F_6$
(the constancy was also confirmed at $8$ random rational points). For the
fibers, Proposition~\ref{prop:curvefib} gives the values of $w_2$ and $w_3$
along the fiber in closed form: writing them as $p_1(w_1)$ and $p_2(w_1)$,
\[
p_2(w_1)=Y_4-3Y_2w_1^2+2Y_1w_1^3-w_1^4,\qquad
p_1(w_1)=Y_3-2Y_2w_1+Y_1w_1^2-w_1^3-w_1^2\,p_2(w_1),
\]
\[
R_Y(w_1)\;=\;p_1(w_1)^2+p_1(w_1)+w_1\,p_2(w_1)+w_1^3+Y_1w_1-Y_2 ,
\]
of degree exactly $12$ with constant leading coefficient $1$ (the identity
$R_Y=$ residual of the second sweep equation was checked at $150$ random
$(Y,w_1)$). At the rational witness target $Y=(1,2,-1,3)$ the polynomial
$R_Y$ is squarefree and coprime to $\gamma(w_1)=Y_1-X_1(w_1,p_1(w_1),p_2(w_1))$
--- both certified by exact univariate gcd computations --- so that fiber has
exactly $12$ points, and by openness of both conditions the generic count is
$12$. Over $Y=0$,
\[
R_0(w_1)\;=\;w_1^{\,5}\,(w_1-1)\,
\bigl(w_1^6+w_1^5+w_1^4-w_1^3-w_1^2-w_1+1\bigr),
\]
the quintuple root $w_1=0$ has $\gamma=0$ and is deleted by the twist, while
$\gamma(1)=1$ and $\gcd(\gamma,R_0/w_1^5)=1$ exactly, so all seven remaining
roots survive: the fiber over $(C_0,0,0,0,0)$ has exactly $7$ points,
including the visible fixed point $(C_0,0,0,0,0)$.
\end{proof}

The first component is
\[
F_{7,1}\;=\;x-61x^2y+9012x^3y^2-\tfrac{35823670}{19}x^4y^3
+\tfrac{238}{19}x^4yz_1+8x^4z_2+x^5z_3 ,
\]
and the remaining components have between $14{,}000$ and $26{,}000$ terms;
they are reproduced, with all verification scripts, in the ancillary files.
Together, $F_6$ and $F_7$ illustrate the two mechanisms available in dimension
five: a mixed direction field on its flexible bottom branch, and the
curve-type field with transverse Keller-type data.

\section{Concluding remarks}\label{sec:conclusion}

The six maps $F$, $G$, $F_4$, $F_5$, $F_6$, $F_7$ studied here are instances of one principle: a
classical enveloping construction (tangent lines of a curve, a tangent direction
field on a surface) supplies an unramified-off-a-divisor, many-to-one polynomial
map whose Jacobian is essentially a coordinate; a monomial twist then launders
that coordinate into the constant Jacobian, at the cost of deleting the
ramification divisor --- so the colliding sheets escape to infinity rather than
merging. The fiber stratifications are correspondingly classical: contact orders
of tangent objects through the target point. For $F$ this gives fiber sizes
$\{3,1,0\}$ (Theorem~\ref{thm:strat3}); for $G$, Theorem~\ref{thm:fibG} proves
the sizes are exactly $\{4,3,2,1,0\}$ --- the value $2$ from double contact at
smooth curve points, $1$ at the cusps, $0$ at the node, and, unexpectedly, $3$
over the hyperplane where the shadow map degenerates; for $F_4$ and $F_5$ the
discriminantal geometry of the swept families is two-dimensional and the
complete stratifications are open, while for $F_6$ and $F_7$ the univariate
fiber polynomials $R_Y$ from the proofs of Theorems~\ref{thm:F6}
and~\ref{thm:F7} reduce the
stratifications to discriminants of a sextic and a degree-twelve polynomial in
the target coordinates, which we have not yet unwound. We will
return to the stratifications of $F_4$, $F_5$, $F_6$ and $F_7$, and to certified
Gr\"obner bases for their graph ideals in favourable orders, in a subsequent
version of this paper. The $M$-branch of \S\ref{subsec:specII} is largely
unexplored: every polynomial solution $(\rho,\lambda)$ of the linear equation
\eqref{eq:wave} is a candidate surface, and the fiber geometry of the resulting
degree-ten (and higher) maps remains to be understood.

Several questions suggest themselves. Which fiber-size sets are realizable by
Keller maps in dimension $n$? (For $F$ no fiber has size one less than the
generic size, yet $G$ does attain size $4-1=3$ --- over the degenerate hyperplane
rather than over a contact stratum --- so the gap phenomenon is not universal and
seems to be contact-geometric in nature.) Do
non-graph surfaces in $\C^4$ admit hyperbolic sweep normalizations, which would
give four-dimensional examples with two independent escape divisors? Is the
middle branch of the mixed direction field $(1,w_1,w_1^2,w_2)^{\mathsf T}$ genuinely
rigid --- that is, does the master equation \eqref{eq:master5} admit any
polynomial solution that is not a suspension of a four-dimensional sweep
(Theorem~\ref{thm:collapse5} and Problem~\ref{prob:rigid5})? More generally, for
which pairs (direction field, branch index) in dimension $n$ is the branch
realizable by irreducibly $n$-dimensional data? And what is
the minimal geometric degree of a Keller counterexample in dimension $n\ge4$?

\subsection*{Acknowledgements}
The computations reported here were carried out in exact rational arithmetic;
scripts are available from the author.


\appendix

\section{Graph ideals, Gr\"obner bases, and exact fiber sizes}\label{app:fibers}

For a polynomial map $\Phi=(\varphi_1,\dots,\varphi_n)^{\mathsf T}\colon\C^n\to\C^n$, the
\emph{graph ideal} is
\[
J_\Phi\;=\;\bigl(\varphi_1-v_1,\dots,\varphi_n-v_n\bigr)\;\subseteq\;
\Q[x_1,\dots,x_n,v_1,\dots,v_n],
\]
the kernel of the substitution homomorphism fixing the $x_i$ and sending
$v_i\mapsto\varphi_i$. Specializing $v$ at a point $v^0$ maps $J_\Phi$ into the
fiber ideal $(\varphi_1-v_1^0,\dots,\varphi_n-v_n^0)$, so any element of
$J_\Phi$ yields constraints valid on every fiber; and since our maps are \'etale,
all fibers are reduced, so ``fiber size'' is unambiguous. All computations below
are in exact rational arithmetic and were certified as indicated.

\subsection{The map $F$}\label{app:F}

For Alp\"oge's map $F$ of \S\ref{sec:first} the reduced lexicographic Gr\"obner
basis of $J_F$ with $z>y>x>v_1>v_2>v_3$ is remarkably small: it consists of six
polynomials $h_1,\dots,h_6$ (integer-scaled below; divide by the leading
coefficients $27,54,6,24,4,8$ for the monic reduced basis). Terms are grouped
by the monomials in $(z,y,x)$, in decreasing lexicographic order, so that each
coefficient is a polynomial in $v=(v_1,v_2,v_3)$:
{\small
\begin{align*}
h_1 ={}& \bigl(27v_1^2v_3^2-18v_1v_2v_3+16v_1+v_2^3v_3-v_2^2\bigr)x^3+\bigl(-3v_2v_3+4\bigr)x-2v_3\\[2pt]
h_2 ={}& \bigl(54v_1v_3^2-18v_2v_3+16\bigr)y\\
&+\bigl(729v_1^3v_3^3-567v_1^2v_2v_3^2+432v_1^2v_3+27v_1v_2^3v_3^2+27v_1v_2^2v_3-48v_1v_2-3v_2^4v_3+3v_2^3\bigr)x^2\\
&+\bigl(-162v_1^2v_3^2+108v_1v_2v_3-96v_1-6v_2^3v_3+6v_2^2\bigr)x\\
&-81v_1v_2v_3^2+72v_1v_3+15v_2^2v_3-16v_2\\[2pt]
h_3 ={}& \bigl(6v_2v_3-8\bigr)yx+6v_3y+\bigl(81v_1^2v_3^2-54v_1v_2v_3+48v_1+3v_2^3v_3-3v_2^2\bigr)x^2\\
&+\bigl(-18v_1v_3+2v_2\bigr)x-3v_2v_3\\[2pt]
h_4 ={}& \bigl(24v_1-2v_2^2\bigr)yx+\bigl(-18v_1v_3+4v_2\bigr)y\\
&+\bigl(-243v_1^3v_3^2+162v_1^2v_2v_3-144v_1^2-9v_1v_2^3v_3+9v_1v_2^2\bigr)x^2\\
&+\bigl(54v_1^2v_3-30v_1v_2+2v_2^3\bigr)x+27v_1v_2v_3-4v_2^2\\[2pt]
h_5 ={}& 4y^2+\bigl(-18v_1v_3-2v_2\bigr)y\\
&+\bigl(-243v_1^3v_3^2+162v_1^2v_2v_3-144v_1^2-9v_1v_2^3v_3+9v_1v_2^2\bigr)x^2\\
&+27v_1v_2v_3-2v_2^2\\[2pt]
h_6 ={}& 8z+\bigl(-108v_1v_3+12v_2\bigr)y\\
&+\bigl(-1458v_1^3v_3^2+243v_1^2v_2^2v_3^2+972v_1^2v_2v_3-864v_1^2-216v_1v_2^3v_3+198v_1v_2^2+9v_2^5v_3-9v_2^4\bigr)x^2\\
&+\bigl(-162v_1^2v_2v_3^2+108v_1v_2^2v_3-96v_1v_2-6v_2^4v_3+6v_2^3\bigr)x\\
&+108v_1^2v_3^2+90v_1v_2v_3-8v_1-23v_2^3v_3+20v_2^2
\end{align*}
}
The basis was certified independently of the algorithm that produced it: all
fifteen $S$-polynomials reduce to zero modulo the set; each $h_i$ vanishes
identically under $v_j\mapsto f_j$; each $f_i-v_i$ reduces to zero modulo the
set; and an independent computation over $\mathbb{F}_p$, $p=2147483629$,
reproduces the basis modulo $p$. The element $h_1$ generates the elimination
ideal $J_F\cap\Q[x,v]$: written by its $(z,y,x)$-grouping,
$h_1=c_3x^3+(4-3v_2v_3)x-2v_3$, and this is the elimination cubic referred to
in Theorem~\ref{thm:strat3}(1); the identities
$\operatorname{disc}_x(h_1)=-c_3a_2^2$ and $a_2^2-4(4-3v_2v_3)^3=108v_3^2c_3$
then drive the whole stratification, whose proof is completed by three lifting
certificates: exact identities $F(\text{candidate point})=v$ verified in
$\Q(v)[x]/(h_1)$, over the function field of $\{a_2=0\}$, and in
$\Q(v_2,v_3)[v_1]/(c_3)$ respectively. This yields
Theorem~\ref{thm:strat3}: fiber sizes $3,1,0$ across the strata, and never
$2$.

\subsection{The map $G$}\label{app:G}

For the degree-four map $G$ of \S\ref{sec:deg4} the behaviour of the
lexicographic basis changes dramatically. The computation, first abandoned
after modular experiments showed staircase elements of $v_1$-degree exceeding
$20$, has since been completed (Singular via Sage): the reduced lexicographic
basis of $J_G$, in the order $x>y>z>v_1>v_2>v_3$, has $174$ elements ---
against six for $F$ --- an interesting contrast, and a warning that the
tangent-sweep counterexamples need not have small Gr\"obner bases. Its profile:
all elements are monic; $20$ leading terms are $x$-headed of $x$-degree
$\le2$ (among them the reconstruction relations with leading terms $xy^2$ and
$xz$), $153$ are $y$-headed of $y$-degree $\le3$, and exactly one element is
free of both $x$ and $y$ --- the elimination quartic $Z_G$ below. Total
degrees range from $12$ to $93$, with exactly one element in each degree from
$61$ to $93$: a staircase ladder of $33$ elements of up to $2{,}787$ terms
(the basis occupies $139$\,MB in plain text), which is where the computation
spends its effort. The fiber analysis, however, does not require the full
basis. Write
$\sigma(v)=(X,Y)^{\mathsf T}:=(v_1v_2,\,v_1^2v_3)^{\mathsf T}$ for the shadow map ($X,Y$ are the values of
$p(w)+2\gamma$ and $q(w)+\gamma w$). The elimination ideal $J_G\cap\Q[x,v]$ is
principal, generated by the quartic (computed as a resultant, and monic in $w$
after the substitution $x=2v_1/(X-p(w))$):
\[
R_G(x;v)\;=\;E(X,Y)\,x^4\;-\;24v_1^2\bigl(3X^2+8X+16Y-24\bigr)x^2
\;-\;64v_1^3(X+4)\,x\;-\;16v_1^4 ,
\]
where
\[
E(X,Y)=27X^4+80X^3-672X^2Y+1536XY^2-512Y^3-144X^2-2304XY+4608Y^2+3456Y
\]
is the implicit equation of the swept quartic curve $\mathcal{K}_G=(p,q)$. As for
$F$, the leading coefficient is the curve equation:
$\operatorname{Res}_w\bigl(W,\,X-p(w)\bigr)=-\tfrac{1}{64}E$ and
$\operatorname{disc}_w(W)=-\tfrac{1}{64}E$ for the tangency quartic
$W=\tfrac14w^4-2w^3+3w^2-Xw+2Y$. The curve $\mathcal{K}_G$ has its node at
$(X,Y)=(-4,\tfrac12)$ --- where the tangency quartic becomes the perfect square
$W=\tfrac14(w^2-4w-2)^2$ --- and its two cusps at the parameters $w=2\pm\sqrt2$,
i.e.\ at $(X,Y)=(-4\mp4\sqrt2,\,-\tfrac{11}2\mp4\sqrt2)$.

The full basis completes this picture from the target side. Its unique
element free of $x$ and $y$ generates the elimination ideal
$J_G\cap\Q[z,v]$: a quartic in $z$ ($114$ terms, coefficients of degree up to
$21$ in $v$),
\[
Z_G(z;v)\;=\;v_1^{\,2}\,z^4\;+\;c_3(v)\,z^3\;+\;c_2(v)\,z^2
\;+\;c_1(v)\,z\;+\;c_0(v),
\]
\[
c_3(v)=-10v_1^3v_2^3+288v_1^3v_2v_3-450v_1^2v_2^2+1504v_1^2v_3
-3648v_1v_2-8208 .
\]
Three exact certificates tie $Z_G$ to the fiber theory below. First,
substituting the reconstruction formula $z=\gamma(6\gamma-4w-\gamma^2)/v_1^2$
of Theorem~\ref{thm:fibG} gives $Z_G\bigl(z(w);v\bigr)\equiv0 \pmod W$
identically (verified in exact arithmetic at random rational targets), so the
four roots of $Z_G$ are precisely the values of $z$ on the four sheets ---
the $z$-coordinate analogue of $R_G$, of degree equal to the geometric
degree. Second, the leading coefficient is exactly $v_1^{\,2}$, while
$c_3(0,v_2,v_3)=-8208\neq0$: over the hyperplane $v_1=0$ the quartic drops to
degree exactly $3$, matching the three-point fibers of
Theorem~\ref{thm:fibG} --- where $R_G$ sees the non-properness locus through
its leading coefficient $E(\sigma(v))$ (the curve side), $Z_G$ sees it
through $v_1^{\,2}$ (the hyperplane side). Third, the discriminant of $Z_G$
with respect to $z$ vanishes at targets with $\sigma(v)\in\mathcal K_G$ and
not at generic targets (verified exactly at sample points of each stratum):
the contact strata of the dual curve are read off the collisions of the
$z$-roots.

\begin{theorem}[Fiber sizes of $G$]\label{thm:fibG}
Every fiber of $G$ is reduced, and its cardinality is determined as follows.
For $v_1\ne0$, fiber points correspond bijectively to roots $w$ of $W$ with
$\gamma:=\tfrac12(X-p(w))\ne0$, via
\[
x=\frac{v_1}{\gamma},\qquad y=\frac{w-\gamma}{v_1},\qquad
z=\frac{\gamma\,(6\gamma-4w-\gamma^2)}{v_1^2},
\]
and these formulas satisfy $G(x,y,z)=v$ identically modulo $W$. Hence for
$v_1\neq0$ the fiber has
\[
\begin{cases}
4\ \text{points} & \text{if }\sigma(v)\notin\mathcal{K}_G,\\
2\ \text{points} & \text{if }\sigma(v)\in\mathcal K_G\text{ is a smooth, non-node point},\\
1\ \text{point} & \text{if }\sigma(v)\text{ is a cusp},\\
0\ \text{points} & \text{if }\sigma(v)\text{ is the node}.
\end{cases}
\]
For $v_1=0$ the quartic $R_G$ degenerates identically and the fiber splits into
the branch $x=0$, which contributes the single point
$\bigl(0,\tfrac{v_2}2,z^*\bigr)$ with $z^*$ determined linearly, and the branch
$\gamma=0$, which contributes the two points with $x^2=4/(v_2^2-24v_3)$; hence
the fiber has $3$ points if $v_2^2\ne24v_3$ and $1$ point if $v_2^2=24v_3$. In
particular $G$ attains every fiber size in $\{4,3,2,1,0\}$.
\end{theorem}

All counts were verified exactly at rational (or explicitly algebraic) sample
points of each stratum, including the exact fiber point
$(0,\tfrac12,-\tfrac14)$ over $v=(0,1,1)$ and the identity
$W=\tfrac14(w^2-4w-2)^2$ at the node. The completed Gr\"obner basis provides
an independent certification: all $174$ elements vanish on the graph of $G$
(checked in modular arithmetic at random rational points of the graph), the
leading terms are pairwise non-divisible, and the quartic $Z_G$ re-derives
the counts over both strata of the target. Note the mechanism of the empty fiber: at
the node \emph{both} branches of the curve impose a double contact, so all four
sheets escape simultaneously --- the analogue for $G$ of the cusp mechanism in
$F$. The value $3$, impossible for $F$, appears over the hyperplane $v_1=0$,
where the shadow map degenerates.

\subsection{The maps $F_4$--$F_7$}\label{app:F4}

For the higher-dimensional maps the graph ideals live in $8$ or $10$
variables, with generators of degree up to $21$, $16$, $44$ and $95$
respectively, and their Gr\"obner bases are out of reach of our implementation
--- the growth from six elements for $F$ to $174$ for $G$ suggests they may be
out of reach altogether. The fiber theory, however, does not need them: for
$F_4$ the fibers are governed by the tangency quintic displayed in the proof
of Theorem~\ref{thm:F4}, whose leading coefficient is the \emph{constant}
$\tfrac29$ (so the number of tangency parameters never drops and all
degenerations pass through $\gamma=0$); for $F_5$ by a resultant of degree
$10$ (proof of Theorem~\ref{thm:F5}); and for $F_6$ and $F_7$ by the
univariate fiber polynomials $R_Y$, of degrees $6$ and $12$ with constant
leading coefficients, derived in closed form in the proofs of
Theorems~\ref{thm:F6} and~\ref{thm:F7}. All four generic counts are
established there by exact elimination and gcd certificates at rational
targets. The complete stratifications --- for $F_4$ a two-dimensional
family of plane curves, for $F_6$ and $F_7$ the discriminants of $R_Y$ ---
are left to a subsequent version.

\begin{thebibliography}{99}

\bibitem{BCW1982}
H.~Bass, E.~H.~Connell and D.~Wright,
\emph{The Jacobian conjecture: reduction of degree and formal expansion of the
inverse}, Bull. Amer. Math. Soc. (N.S.) \textbf{7} (1982), 287--330.

\bibitem{BBR1962}
A.~Bia{\l}ynicki-Birula and M.~Rosenlicht,
\emph{Injective morphisms of real algebraic varieties},
Proc. Amer. Math. Soc. \textbf{13} (1962), 200--203.

\bibitem{CourantHilbert1962}
R.~Courant and D.~Hilbert,
\emph{Methods of Mathematical Physics. Vol.~II: Partial Differential
Equations}, Interscience, New York, 1962.

\bibitem{Druzkowski1983}
L.~M.~Dru\.zkowski,
\emph{An effective approach to Keller's Jacobian conjecture},
Math. Ann. \textbf{264} (1983), 303--313.

\bibitem{Essen2000}
A.~van den Essen,
\emph{Polynomial Automorphisms and the Jacobian Conjecture},
Progress in Mathematics \textbf{190}, Birkh\"auser, Basel, 2000.

\bibitem{Jelonek1993}
Z.~Jelonek,
\emph{The set of points at which a polynomial map is not proper},
Ann. Polon. Math. \textbf{58} (1993), 259--266.

\bibitem{Keller1939}
O.-H.~Keller,
\emph{Ganze Cremona-Transformationen},
Monatsh. Math. Phys. \textbf{47} (1939), 299--306.

\bibitem{Moh1983}
T.-T.~Moh,
\emph{On the Jacobian conjecture and the configurations of roots},
J. Reine Angew. Math. \textbf{340} (1983), 140--212.

\bibitem{Pinchuk1994}
S.~Pinchuk,
\emph{A counterexample to the strong real Jacobian conjecture},
Math. Z. \textbf{217} (1994), 1--4.

\bibitem{RosenbloomWidder1959}
P.~C.~Rosenbloom and D.~V.~Widder,
\emph{Expansions in terms of heat polynomials and associated functions},
Trans. Amer. Math. Soc. \textbf{92} (1959), 220--243.

\bibitem{Alpoge2026}
L.~Alp\"oge,
\emph{A counterexample to the Jacobian conjecture in dimension three},
announcement, July 19, 2026,
\url{https://x.com/__alpoge__/status/2079028340955197566}; see also
\cite{SBS2026,Tao2026}.
% TODO: replace by the formal reference when available.

\bibitem{Gallagher2026}
A.~Gallagher,
\emph{An infinite family of counterexamples to the Jacobian Conjecture in
dimension three: every generic fiber degree $n\ge3$ occurs},
preprint, July 20, 2026, Zenodo, \href{https://doi.org/10.5281/zenodo.21479195}{doi:10.5281/zenodo.21479195}.

\bibitem{SBS2026}
Secret Blogging Seminar,
\emph{The new counterexample to the Jacobian conjecture}, blog post, July 20,
2026,
\url{https://sbseminar.wordpress.com/2026/07/20/the-new-counterexample-to-the-jacobian-conjecture/}.

\bibitem{Speyer2026}
D.~E.~Speyer,
\emph{The geometry and structure of Gallagher's counterexamples to the Jacobian
conjecture}, note, July 23, 2026,
\url{https://sbseminar.wordpress.com/wp-content/uploads/2026/07/jacobiantangentsweep.pdf}.

\bibitem{Tao2026}
T.~Tao,
\emph{A digestion of the Jacobian conjecture counterexample}, blog post, July 21,
2026,
\url{https://terrytao.wordpress.com/2026/07/21/a-digestion-of-the-jacobian-conjecture-counterexample/}.

\bibitem{Wang1980}
S.~S.-S.~Wang,
\emph{A Jacobian criterion for separability},
J. Algebra \textbf{65} (1980), 453--494.

\bibitem{Yagzhev1980}
A.~V.~Yagzhev,
\emph{On Keller's problem},
Sibirsk. Mat. Zh. \textbf{21} (1980), 141--150.

\end{thebibliography}

\end{document}

